%% =========================================================
%% Chapter: Physics examples of the multi-block asymmetric
%% compression theorem.
%% Covers Notes/OpenProblemsTN/strategies/p5_asymmetric_compression_theorem.tex,
%% Examples A--E, the further physics examples of
%% Notes/OpenProblemsTN/checks/p5_more_examples_data.md, and the
%% renormalization fixed points of
%% Notes/OpenProblemsTN/checks/p6_examples_compression_data.md.
%% =========================================================

\chapter{Physics Examples of the Multi-Block Compression Theorem}%
\label{ch:asymmetric_examples}

Matrix product states and operators describe periodic spin chains by
contracting local tensors around a ring~\cite{Cirac2021Matrix}. At every
positive chain length $N$, an MPS tensor $A$ and an MPO tensor $M$ give
\begin{align}
    \ket{V^{(N)}(A)}
     & =\sum_{i_1,\ldots,i_N}\tr(A^{i_1}\cdots A^{i_N})
    \ket{i_1,\ldots,i_N}, \notag                        \\
    O_N(M)
     & =\sum_{i_1,\ldots,i_N}\sum_{j_1,\ldots,j_N}
    \tr(M^{i_1j_1}\cdots M^{i_Nj_N})
    \ket{i_1,\ldots,i_N}\bra{j_1,\ldots,j_N}.
    \notag
\end{align}
The physical indices $i,j$ range over $\{0,\ldots,d-1\}$.
Each local matrix acts on its tensor's virtual bond space, with rows
and columns indexed by $\{0,\ldots,D-1\}$ for bond dimension $D$.
The periodic boundary contraction is the ordinary matrix trace, with
no boundary insertion. The first MPO physical index is the output and
the second is the input; operator products act from right to left.
Thus the local tensor $\sum_jM^{ij}\otimes P^{jk}$ represents
$O_N(M)O_N(P)$. MPS vectors need not be normalized. A nonzero positive
periodic MPDO operator gives a density operator after division by its
trace.

Physical examples of matrix product operators, topological symmetries, and
fixed points illustrate the multi-block asymmetric compression theorem
(Theorem~\ref{thm:asym_multiblock}).


\section{Explicit gauges and exact arithmetic}%
\label{sec:asymex_explicit_gauge}

The compression datum of Theorem~\ref{thm:asym_multiblock} carries its
change of bond coordinates as an invertible map onto a graded coordinate
space, one summand per slot. For each worked example below, that map is an explicit
invertible matrix together with an explicit labelling of the bond
coordinates by slots and positions inside slots.

\begin{definition}[Explicit gauge and entrywise integer coercion]
    \label{def:asymex_explicit_gauge}
    \lean{MPSTensor.complexOfInt, MPSTensor.mulIntTensor, MPSTensor.gaugeOfMatrix}
    \leanok
    \uses{def:mps_tensor, def:mpdo_operator_product}
    Let $\tau$ be a bijection between the graded coordinate space of a slot
    family and the bond index set $\{0,\ldots,D_B-1\}$, and let $G$ and
    $G^{-1}$ be mutually inverse matrices in $\MN{D_B}$. The \emph{explicit
        gauge} attached to $(\tau,G)$ is the invertible map sending a bond
    vector $v$ to the graded vector $b\mapsto(Gv)(\tau b)$. For an integer
    matrix $X$, its \emph{entrywise coercion} is the complex matrix with the
    same entries; for integer matrix product operator tensors $M$ and $N$,
    their integer bond-space product is
    $(M\mathbin{\cdot}N)^{ik}=\sum_jM^{ij}\otimes N^{jk}$.
\end{definition}

\begin{theorem}[Conjugation, coercion and the compression pairs of a slot]
    \label{thm:asymex_explicit_gauge}
    % Principal declaration: MPSTensor.conjMatrix_gaugeOfMatrix.
    \lean{MPSTensor.conjMatrix_gaugeOfMatrix, MPSTensor.complexOfInt_mul,
        MPSTensor.mulTensor_complexOfRing, MPSTensor.mulTensor_smul_complexOfRing,
        MPSTensor.MultiBlockCompression.left_gaugeOfMatrix,
        MPSTensor.MultiBlockCompression.right_gaugeOfMatrix}
    \leanok
    \uses{def:asymex_explicit_gauge, thm:asym_multiblock}
    Conjugation by the explicit gauge attached to $(\tau,G)$ sends a bond
    matrix $A$ to the matrix $GAG^{-1}$ relabelled along $\tau$. The
    entrywise coercion is multiplicative, and the bond-space product of two
    integer tensors each rescaled by one scalar $c$ is $c^2$ times the
    entrywise coercion of their integer bond-space product. If a compression
    datum uses the explicit gauge attached to $(\tau,G)$, then the
    compression map $W_s$ out of the bond space consists of the rows of $G$
    at the coordinates of the slot $s$, and the compression map $V_s$ into
    the bond space consists of the columns of $G^{-1}$ at those same
    coordinates.
\end{theorem}

\begin{proof}
    \leanok
    The matrix of the explicit gauge is $G$ relabelled along $\tau$ in its
    rows, and the matrix of its inverse is $G^{-1}$ relabelled along $\tau$
    in its columns, so conjugation is the stated matrix conjugation followed
    by the relabelling. Multiplicativity of the coercion is the fact that a
    sum of products of integers is the same computed in $\Z$ or in $\C$, and
    the statement about the bond-space product follows by expanding both
    sides entrywise. The two compression maps are the composites of the
    coordinate embedding and projection of the slot with the gauge and its
    inverse; evaluating the matrix products at a single graded coordinate
    leaves one row of $G$ and one column of $G^{-1}$.
\end{proof}

\begin{definition}[One-slot compression data]%
    \label{def:asymex_one_slot}
    \lean{MPSTensor.stackedInt, MPSTensor.oneSlot, MPSTensor.oneSlotMem,
        MPSTensor.MultiBlockCompression.ofConjMatrix,
        MPSTensor.MultiBlockCompression.ofConjInt,
        MPSTensor.MultiBlockCompression.ofEq,
        MPSTensor.MultiBlockCompression.ofScalarFlagFour}
    \leanok
    \uses{def:asymex_explicit_gauge, thm:asym_multiblock}
    For a stacked product of two integer matrix product operator tensors,
    read over the pair alphabet, the letter at $a=do+i$ is
    $(M\mathbin{\cdot}N)^{oi}$. A \emph{one-slot compression datum} for a
    tensor $B$ onto a single target $C$ is the data of
    Theorem~\ref{thm:asym_multiblock} for the one-element slot set: an
    explicit gauge $(\tau,G)$ conjugating $B$ into a matrix that is block
    upper triangular for a flag with $C$ at one diagonal position. Two
    shapes recur in the examples below: a conjugation $GB^iG^{-1}=C^i$ for
    an invertible $G$ gives such a datum with no zero slots ($z=0$); and,
    for a four-dimensional $B$ and a one-dimensional $C$, an integer gauge
    sending every letter to a matrix with $C$ at the third diagonal
    position of a four-step flag and zero at the other three gives such a
    datum with three zero slots ($z=3$).
\end{definition}

\begin{theorem}[The remainder of a conjugate datum vanishes; word traces
        of a one-slot datum]%
    \label{thm:asymex_one_slot}
    \lean{MPSTensor.MultiBlockCompression.remainder_ofConjMatrix,
        MPSTensor.MultiBlockCompression.remainder_ofEq,
        MPSTensor.MultiBlockCompression.remainder_ofConjInt,
        MPSTensor.MultiBlockCompression.trace_evalWord_oneSlot,
        MPSTensor.MultiBlockCompression.trace_evalWord_oneSlot_smul}
    \leanok
    \uses{def:asymex_one_slot}
    The remainder of a conjugate one-slot datum vanishes identically: the
    compression is an exact change of coordinates. For every one-slot
    datum, the trace of every nonempty word of $B$ equals the trace of the
    corresponding word of $C$, or, if the target is carried with a scalar
    weight $c$, that power of $c$ times the trace of the word of the
    unweighted target.
\end{theorem}

\begin{proof}
    \leanok
    \uses{thm:asymex_explicit_gauge, prop:asym_compression_trace}
    Immediate from Theorem~\ref{thm:asymex_explicit_gauge} and
    Proposition~\ref{prop:asym_compression_trace}, specialized to a single
    slot.
\end{proof}

\begin{theorem}[Integer certificates for normality and for vanishing
        sitewise intertwiners]%
    \label{thm:asymex_int_certificates}
    \lean{MPSTensor.isNBlkInjective_two_of_int,
        MPSTensor.right_intertwiner_eq_zero_of_int,
        MPSTensor.left_intertwiner_eq_zero_of_int}
    \leanok
    \uses{def:normal}
    If every matrix unit is, up to a common nonzero integer factor, an
    integer combination of length-two words of an integer tensor, the
    length-two words span the full matrix algebra. If the rows,
    respectively columns, of the letters occurring in a sitewise
    intertwiner system form an integer matrix with an integer left inverse
    up to a nonzero factor, the only solution of that system is zero.
\end{theorem}

\begin{proof}
    \leanok
    Clearing the nonzero integer factor by scalar division after coercion
    to $\C$ reduces each statement to the corresponding identity between
    integer matrices, a finite verification.
\end{proof}

\section{Hidden symmetry-broken sectors}%
\label{sec:asymex_ghz}

The following example describes two broken sectors of a $\Z_2$ symmetry
by a periodic vector family generated by the sum of two inequivalent
one-dimensional normal blocks, one per sector. A tensor generating that same family need
not display the sectors: they may occupy non-adjacent positions, in an order
different from the target's, coupled through a zero block. This is the
smallest example with more than one target slot, and it serves as a warm-up
for the fusion examples below.

\begin{definition}[The two sectors and the four-dimensional source]
    \label{def:asymex_ghz}
    \lean{MPSTensor.ghzC, MPSTensor.ghzB}
    \leanok
    \uses{def:mps_tensor}
    The target is the pair of one-dimensional tensors
    $C_0=\bigl((1),(0)\bigr)$ and $C_1=\bigl((0),(1)\bigr)$ on the alphabet
    $\{0,1\}$, whose direct sum is the Greenberger--Horne--Zeilinger tensor
    $\diag(1,0)$, $\diag(0,1)$. The source is the bond-dimension-four
    tensor
    \begin{align}
        B^0
         & =\begin{pmatrix}
                0 & 1 & 2 & 0 \\
                0 & 0 & 1 & 1 \\
                0 & 0 & 1 & 3 \\
                0 & 0 & 0 & 0
            \end{pmatrix},
         &
        B^1
         & =\begin{pmatrix}
                1 & 1 & 0 & 2 \\
                0 & 0 & 2 & 0 \\
                0 & 0 & 0 & 1 \\
                0 & 0 & 0 & 0
            \end{pmatrix}.
        \notag
    \end{align}
\end{definition}

\begin{center}
    % Formula: the periodic trace tr(B^w) of the bond-dimension-four
    % source of this section.
    % Source: Definition~\ref{def:asymex_ghz}.
    % Ink-to-index map: each node is the source tensor $B$; the upper
    % legs carry the letters of the word.
    % Contraction set: the horizontal virtual legs, closed around the
    % ring by the trace; no leg is left open.
    % Boundary signature: (virtual-west=0 | virtual-east=0 |
    % physical-up=3 | physical-down=0) for the finite schematic shown.
    \tenkzkernel
    \begin{tenkz}[cols=4, boundary=periodic, physical=up]
        \tn[label pos=135, ports={90:physical:$i_1$}]{B} &
        \tn[label pos=135]{B} & \tn[skin=dots]{} &
        \tn[label pos=135, ports={90:physical:$i_N$}]{B}
    \end{tenkz}
\end{center}

\begin{theorem}[Compression of the two sectors]
    \label{thm:asymex_ghz_compression}
    % Principal declaration: MPSTensor.ghzSectors_isReduction.
    \lean{MPSTensor.ghzSectors_compression, MPSTensor.ghzSectors_isReduction,
        MPSTensor.ghzSectors_dim_eq}
    \leanok
    \uses{def:asymex_ghz, thm:asym_multiblock}
    The source of Definition~\ref{def:asymex_ghz} carries the block
    triangular form of Theorem~\ref{thm:asym_multiblock} for the slot family
    $(C_0,C_1)$ with $z=2$ zero slots, the four diagonal positions carrying,
    in order, the sector $C_1$, a zero slot, the sector $C_0$ and a second
    zero slot. The compression pair of each sector satisfies $W_sV_s=\id$
    and $W_sB^wV_s=C_s^w$ for every word, and $4=1+1+2$.
\end{theorem}

\begin{proof}
    \leanok
    \uses{thm:asymex_explicit_gauge}
    The gauge is the relabelling of the four bond coordinates by the four
    slots, with no change of basis. The three matrix conditions of the
    compression datum, block triangularity and the identification of the
    matched and unmatched diagonal blocks, are then a finite verification
    over the integer entries of $B^0$ and $B^1$. The compression pairs and
    the dimension count are clauses (iv), (v) and (vii) of
    Theorem~\ref{thm:asym_multiblock}.
\end{proof}

\begin{theorem}[Periodic data of the four-dimensional source]
    \label{thm:asymex_ghz_trace}
    \lean{MPSTensor.ghzB_trace_evalWord_eq_sum}
    \leanok
    \uses{def:asymex_ghz}
    For every nonempty word $w$, $\tr(B^w)=\tr(C_0^w)+\tr(C_1^w)$.
\end{theorem}

\begin{proof}
    \leanok
    \uses{thm:asymex_ghz_compression, prop:asym_compression_trace}
    This is Proposition~\ref{prop:asym_compression_trace} applied to the
    compression datum of Theorem~\ref{thm:asymex_ghz_compression}, whose
    slot set has two members.
\end{proof}

\begin{theorem}[The first sector has no sitewise intertwiner]
    \label{thm:asymex_ghz_no_intertwiner}
    \lean{MPSTensor.ghzC0_right_intertwiner_eq_zero,
        MPSTensor.ghzC0_left_intertwiner_eq_zero}
    \leanok
    \uses{def:asymex_ghz}
    If $v\in\C^4$ satisfies $B^0v=v$ and $B^1v=0$, then $v=0$; if
    $u\in\C^4$ satisfies $uB^0=u$ and $uB^1=0$, then $u=0$. Both sitewise
    intertwiner spaces for the sector $C_0$ therefore vanish, and the
    word-level compression of Theorem~\ref{thm:asymex_ghz_compression} is
    the strongest local relation available for that sector.
\end{theorem}

\begin{proof}
    \leanok
    The last row of $B^0$ vanishes, so the last coordinate of $B^0v=v$
    gives $v_3=0$; the second coordinate of $B^1v=0$ gives $v_2=0$; and the
    second and first coordinates of $B^0v=v$ then force $v_1=0$ and
    $v_0=0$. For the covector, the first and third coordinates of $uB^1=0$
    give $u_0=0$ and $u_1=0$, its last coordinate gives $u_2=0$, and the
    last coordinate of $uB^0=u$ gives $u_3=0$.
\end{proof}

\section{A repeated block with a nontrivial extension}%
\label{sec:asymex_repeated}

Periodic traces determine the multiplicity with which a simple block occurs,
but not the extension carrying the copies. The following warm-up shows that
the conclusion of the symmetric fundamental theorem, an invertible gauge
relating source and target, is false for such a source, not merely
inapplicable.

\begin{definition}[The scalar block and its Jordan extension]
    \label{def:asymex_repeated}
    \lean{MPSTensor.repA, MPSTensor.repC, MPSTensor.repB}
    \leanok
    \uses{def:mps_tensor}
    The target is the scalar normal tensor $A^0=(1)$, $A^1=(2)$, taken
    twice, and the source is
    \begin{align}
        B^0
         & =\begin{pmatrix}1&1\\0&1\end{pmatrix},
         &
        B^1
         & =\begin{pmatrix}2&0\\0&2\end{pmatrix}.
        \notag
    \end{align}
\end{definition}

\begin{center}
    % Formula: the periodic trace tr(B^w) of the bond-dimension-two
    % Jordan source of this section.
    % Source: Definition~\ref{def:asymex_repeated}.
    % Ink-to-index map: each node is the source tensor $B$; the upper
    % legs carry the letters of the word.
    % Contraction set: the horizontal virtual legs, closed around the
    % ring by the trace; no leg is left open.
    % Boundary signature: (virtual-west=0 | virtual-east=0 |
    % physical-up=3 | physical-down=0) for the finite schematic shown.
    \tenkzkernel
    \begin{tenkz}[cols=4, boundary=periodic, physical=up]
        \tn[label pos=135, ports={90:physical:$i_1$}]{B} &
        \tn[label pos=135]{B} & \tn[skin=dots]{} &
        \tn[label pos=135, ports={90:physical:$i_N$}]{B}
    \end{tenkz}
\end{center}

\begin{theorem}[Compression onto the two copies]
    \label{thm:asymex_repeated_compression}
    % Principal declaration: MPSTensor.repeatedBlock_isReduction.
    \lean{MPSTensor.repeatedBlock_compression, MPSTensor.repeatedBlock_isReduction,
        MPSTensor.repB_trace_evalWord_eq_two_mul}
    \leanok
    \uses{def:asymex_repeated, thm:asym_multiblock}
    The source of Definition~\ref{def:asymex_repeated} carries the block
    triangular form of Theorem~\ref{thm:asym_multiblock} for the two copies
    of $A$ with no zero slots, each coordinate pair compresses its copy,
    $W_sV_s=\id$ and $W_sB^wV_s=A^w$ for every word, and
    $\tr(B^w)=2\tr(A^w)$ for every nonempty word.
\end{theorem}

\begin{proof}
    \leanok
    \uses{prop:asym_compression_trace}
    Both matrices are upper triangular with both diagonal entries equal to
    the corresponding scalar of $A$, which is the whole content of the three
    matrix conditions; the compression pairs are clauses (iv) and (v) of
    Theorem~\ref{thm:asym_multiblock}, and the trace identity is
    Proposition~\ref{prop:asym_compression_trace} with two equal slots.
\end{proof}

\begin{theorem}[No invertible gauge onto the repeated target]
    \label{thm:asymex_repeated_not_gauge}
    \lean{MPSTensor.not_gaugeEquiv_directSum_repA}
    \leanok
    \uses{def:asymex_repeated, def:gauge_equiv}
    There is no invertible $X\in\GL_2(\C)$ with
    $B^i=X\diag(A^i,A^i)X^{-1}$ for both letters $i$.
\end{theorem}

\begin{proof}
    \leanok
    For $i=0$ the matrix $\diag(A^0,A^0)$ is the identity, so any such
    conjugation would give $B^0=\id$, contradicting the nonzero
    off-diagonal entry of $B^0$.
\end{proof}

\section{The square of the CZX $\Z_2$ symmetry}%
\label{sec:asymex_czx}

On a periodic chain of $L$ qubits the symmetry of CZX type is
\begin{align}
    U_L & =\Bigl(\prod_{i=1}^{L}CZ_{i,i+1}\Bigr)X^{\otimes L},
    \label{eq:asymex_czx_operator}
\end{align}
a matrix product operator of bond dimension two with tensor
$M^{ij}=\delta_{i,1\oplus j}T_j$. This is the CZX example of an anomalous
$\Z_2$ symmetry. Its square is $U_L^2=(-1)^L\id^{\otimes L}$, so the stacked product of its
tensor with itself has, at every positive length, the periodic data of the
one-dimensional identity tensor carried with the weight $-1$. The fusion is
therefore beyond the unit-weight setting. The sitewise theory of
\cite[Appendix~A]{GarreRubio2023MPOSym} imposes a closedness requirement
for arbitrary boundary conditions. The compression and intertwiner
calculations below concern this explicit tensor, not a classification of
representations by their three-cocycle.

\begin{definition}[The bond-two symmetry tensor and its square]
    \label{def:asymex_czx_tensors}
    \lean{CZXCompression.czxTensor, CZXCompression.czxMPS,
        CZXCompression.identityTensor, CZXCompression.identityMPS,
        CZXCompression.czxSquare, CZXCompression.czxSquareTarget}
    \leanok
    \uses{def:mpo_tensor, def:mpo_to_mps, def:mpdo_operator_product}
    The symmetry tensor is $M^{ij}=\delta_{i,1\oplus j}T_j$ with
    \begin{align}
        T_0 & =\begin{pmatrix}1&0\\1&0\end{pmatrix},
            &
        T_1 & =\begin{pmatrix}0&1\\0&-1\end{pmatrix};
        \notag
    \end{align}
    the identity tensor is $\delta^{ij}$ of bond dimension one. Read over
    the alphabet of index pairs, the symmetry tensor has the four matrices
    $0$, $T_1$, $T_0$, $0$ in the order $(0,0)$, $(0,1)$, $(1,0)$, $(1,1)$.
    The source is the stacked product
    $B^{ij}=\sum_mM^{im}\otimes M^{mj}$ of bond dimension four, and the
    single target is $-\delta$.
\end{definition}

\begin{center}
    % Formula: B^{ij}=\sum_mM^{im}\otimes M^{mj}, the stacked product of
    % Definition~\ref{def:asymex_czx_tensors}.
    % Source: Definition~\ref{def:asymex_czx_tensors}.
    % Ink-to-index map: both rows are the bond-two symmetry tensor $M$;
    % the shared vertical pairleg at each site is the summed index $m$.
    % Contraction set: the vertical pairlegs (index $m$) at every site
    % and the horizontal virtual bonds within each row; the physical
    % legs $(i,j)$ remain open.
    % Boundary signature: (virtual-west=0 | virtual-east=0 |
    % physical-up=3 | physical-down=3) on both sides.
    \tenkzkernel
    \begin{tenkz}[rows={op,op}, cols=4, boundary=periodic]
        \tn[skin=mpo, ports={90:physical:$i_1$, 270:physical}]{M} &
        \tn[skin=mpo, ports={90:physical:$i_2$, 270:physical}]{M} &
        \tn[skin=dots, ports={180:virtual, 0:virtual}]{} &
        \tn[skin=mpo, ports={90:physical:$i_N$, 270:physical}]{M} \\
        \tn[skin=mpo, ports={90:physical, 270:physical:$j_1$}]{M} &
        \tn[skin=mpo, ports={90:physical, 270:physical:$j_2$}]{M} &
        \tn[skin=dots, ports={180:virtual, 0:virtual}]{} &
        \tn[skin=mpo, ports={90:physical, 270:physical:$j_N$}]{M}
    \end{tenkz}
    \quad $=$ \quad
    \begin{tenkz}[rows={op}, cols=4, boundary=periodic]
        \tn[skin=mpo, ports={90:physical:$i_1$, 270:physical:$j_1$}]{B} &
        \tn[skin=mpo, ports={90:physical:$i_2$, 270:physical:$j_2$}]{B} &
        \tn[skin=dots, ports={180:virtual, 0:virtual}]{} &
        \tn[skin=mpo, ports={90:physical:$i_N$, 270:physical:$j_N$}]{B}
    \end{tenkz}
\end{center}

\begin{theorem}[Normality of the symmetry tensor and of the target]
    \label{thm:asymex_czx_normal}
    \lean{CZXCompression.czxMPS_isNormal, CZXCompression.czxSquareTarget_isNormal}
    \leanok
    \uses{def:asymex_czx_tensors, def:normal}
    The bond-two symmetry tensor is normal: its length-two words span
    $\MN{2}$. The target $-\delta$ is normal: its bond dimension is one and
    one of its matrices is the nonzero scalar $-1$.
\end{theorem}

\begin{proof}
    \leanok
    \uses{thm:asymex_normal_units}
    Each of the four matrix units of $\MN{2}$ is a half-sum or a
    half-difference of two of the four products $T_aT_b$, which is a finite
    verification over integer matrices; the span of the length-two words is
    therefore all of $\MN{2}$. For the target, its matrix at the letter
    $(0,0)$ is the nonzero scalar $-1$, so
    Theorem~\ref{thm:asymex_normal_units} applies with the single position
    of the one-by-one matrix algebra.
\end{proof}

\begin{theorem}[Compression of the squared symmetry]
    \label{thm:asymex_czx_compression}
    % Principal declaration: CZXCompression.czxSquare_isReduction.
    \lean{CZXCompression.czxSquare_compression, CZXCompression.czxSquareLeft,
        CZXCompression.czxSquareRight, CZXCompression.czxSquare_isReduction,
        CZXCompression.czxSquareLeft_mul_right, CZXCompression.czxSquare_dim_eq,
        CZXCompression.czxSquare_evalWord_remainder_eq_zero}
    \leanok
    \uses{def:asymex_czx_tensors, thm:asym_multiblock}
    The stacked product of Definition~\ref{def:asymex_czx_tensors} carries
    the block triangular form of Theorem~\ref{thm:asym_multiblock} for the
    single slot $-\delta$ with $z=3$ zero slots and $4=1+3$. The
    compression pair is
    \begin{align}
        W & =(0,0,-1,0),
          &
        V & =(1,1,-1,-1)^{\mathsf T},
        \notag
    \end{align}
    with $WV=1$ and $WB^wV=(-\delta)^w$ for every word, and the remainder
    $R^{ij}=B^{ij}-V(-\delta)^{ij}W$ satisfies
    $R^{i_1j_1}R^{i_2j_2}R^{i_3j_3}R^{i_4j_4}=0$.
\end{theorem}

\begin{proof}
    \leanok
    \uses{thm:asymex_explicit_gauge}
    The change of bond coordinates has as its columns, in the order of the
    flag, the line $(1,0,0,-1)$ annihilated by every letter, the vector
    $(1,1,-1,-1)$ spanning the image of every letter modulo that line, and
    two arbitrary completions. In these coordinates every letter is block
    upper triangular with the entry $-1$ in the second diagonal position and
    zero in the other three, which is a finite verification over integer
    matrices. The explicit pair is
    Theorem~\ref{thm:asymex_explicit_gauge} applied to the second
    coordinate, and the vanishing of a product of four remainder matrices is
    clause (vi) of Theorem~\ref{thm:asym_multiblock} with $|S|+z=4$.
\end{proof}

\begin{theorem}[Periodic data of the squared symmetry]
    \label{thm:asymex_czx_trace}
    \lean{CZXCompression.czxSquare_trace_evalWord}
    \leanok
    \uses{def:asymex_czx_tensors}
    For every nonempty word $w$ over the alphabet of index pairs,
    $\tr(B^w)=(-1)^{|w|}\tr(\delta^w)$: this is
    $U_L^2=(-1)^L\id^{\otimes L}$ at the level of periodic coefficients.
\end{theorem}

\begin{proof}
    \leanok
    \uses{thm:asymex_czx_compression, prop:asym_compression_trace}
    Proposition~\ref{prop:asym_compression_trace} applied to the datum of
    Theorem~\ref{thm:asymex_czx_compression} gives
    $\tr(B^w)=\tr\bigl((-\delta)^w\bigr)$, and pulling the scalar out of the
    word evaluation contributes the factor $(-1)^{|w|}$.
\end{proof}

\begin{theorem}[Both sitewise intertwiner spaces vanish]
    \label{thm:asymex_czx_no_intertwiner}
    \lean{CZXCompression.czxSquare_right_intertwiner_eq_zero,
        CZXCompression.czxSquare_left_intertwiner_eq_zero}
    \leanok
    \uses{def:asymex_czx_tensors}
    If $v\in\C^4$ satisfies $B^{ij}v=(-\delta)^{ij}v$ for every index pair,
    then $v=0$; if $u\in\C^4$ satisfies $uB^{ij}=(-\delta)^{ij}u$ for every
    index pair, then $u=0$.
\end{theorem}

\begin{proof}
    \leanok
    Only the letters $(0,0)$ and $(1,1)$ of the stacked tensor are nonzero,
    and each is a rank-one matrix with an explicit integer kernel and image.
    Comparing the four coordinates of the two equations for those two
    letters gives, in turn, $v_1=v_2=v_0=v_3=0$, and the transposed
    comparison gives $u_0=u_1=u_3=u_2=0$.
\end{proof}

\begin{remark}[Failure of closure under conjugate transposition]
    \label{rem:asymex_czx_anomaly}
    \uses{cor:asym_anomaly_not_star}
    Theorem~\ref{thm:asymex_czx_no_intertwiner} is what
    Corollary~\ref{cor:asym_anomaly_not_star} converts into the statement
    that the algebra generated by the matrices of the stacked tensor is not
    closed under conjugate transposition, although the symmetry operator
    itself is unitary. The obstruction to a sitewise fusion pair is
    therefore algebraic and not a defect of the particular coordinates
    chosen here. This algebraic implication does not identify nonsplitting
    with an anomaly class. Determining that class requires a separate
    computation of the fusion-tensor associator and its cohomology class.
\end{remark}

\section{A non-simple symmetry family with an oscillating coefficient}%
\label{sec:asymex_czx_plus}

Adding the identity to the symmetry of Section~\ref{sec:asymex_czx} gives
the non-simple family $O_L=\id^{\otimes L}+U_L$, generated by the
bond-dimension-three tensor $N=\delta\oplus M$ whose two blocks are normal
and inequivalent. Its square is
\begin{align}
    O_L^2 & =\bigl(1+(-1)^L\bigr)\id^{\otimes L}+2\,U_L,
    \label{eq:asymex_czx_plus_square}
\end{align}
an operator family whose identity structure coefficient oscillates with the length
while remaining an integer. The compression resolves the oscillation
exactly as Theorem~\ref{thm:asym_power_sum_coefficients} predicts: the
identity block occurs twice, with the weights $+1$ and $-1$, and
$1+(-1)^L$ is the length-$L$ power sum of the weight multiset
$\{+1,-1\}$. This is an example of the general power-sum mechanism,
not of the positive diagonal matrices in the CPSV characterization
of MPDO renormalization fixed points.

\begin{definition}[The non-simple family and its square]
    \label{def:asymex_czx_plus}
    \lean{CZXCompression.plusTensor, CZXCompression.czxPlusIdentity,
        CZXCompression.plusBlockDim, CZXCompression.plusTargets}
    \leanok
    \uses{def:asymex_czx_tensors, def:mpdo_operator_product}
    The generating tensor is $N^{ij}=\delta^{ij}\oplus M^{ij}$ of bond
    dimension three, the source is the stacked product
    $B^{ij}=\sum_mN^{im}\otimes N^{mj}$ of bond dimension nine, and the four
    target slots are $\delta$, $-\delta$, $M$ and $M$, of bond dimensions
    $1$, $1$, $2$ and $2$.
\end{definition}

\begin{center}
    % Formula: B^{ij}=\sum_mN^{im}\otimes N^{mj}, the stacked product of
    % Definition~\ref{def:asymex_czx_plus}.
    % Source: Definition~\ref{def:asymex_czx_plus}.
    % Ink-to-index map: both rows are the bond-three generating tensor
    % $N=\delta\oplus M$; the shared vertical pairleg at each site is
    % the summed index $m$.
    % Contraction set: the vertical pairlegs (index $m$) at every site
    % and the horizontal virtual bonds within each row; the physical
    % legs $(i,j)$ remain open.
    % Boundary signature: (virtual-west=0 | virtual-east=0 |
    % physical-up=3 | physical-down=3) on both sides.
    \tenkzkernel
    \begin{tenkz}[rows={op,op}, cols=4, boundary=periodic]
        \tn[skin=mpo, ports={90:physical:$i_1$, 270:physical}]{N} &
        \tn[skin=mpo, ports={90:physical:$i_2$, 270:physical}]{N} &
        \tn[skin=dots, ports={180:virtual, 0:virtual}]{} &
        \tn[skin=mpo, ports={90:physical:$i_N$, 270:physical}]{N} \\
        \tn[skin=mpo, ports={90:physical, 270:physical:$j_1$}]{N} &
        \tn[skin=mpo, ports={90:physical, 270:physical:$j_2$}]{N} &
        \tn[skin=dots, ports={180:virtual, 0:virtual}]{} &
        \tn[skin=mpo, ports={90:physical, 270:physical:$j_N$}]{N}
    \end{tenkz}
    \quad $=$ \quad
    \begin{tenkz}[rows={op}, cols=4, boundary=periodic]
        \tn[skin=mpo, ports={90:physical:$i_1$, 270:physical:$j_1$}]{B} &
        \tn[skin=mpo, ports={90:physical:$i_2$, 270:physical:$j_2$}]{B} &
        \tn[skin=dots, ports={180:virtual, 0:virtual}]{} &
        \tn[skin=mpo, ports={90:physical:$i_N$, 270:physical:$j_N$}]{B}
    \end{tenkz}
\end{center}

\begin{theorem}[Compression of the squared non-simple family]
    \label{thm:asymex_czx_plus_compression}
    % Principal declaration: CZXCompression.czxPlusIdentity_isReduction.
    \lean{CZXCompression.czxPlusIdentity_compression,
        CZXCompression.czxPlusIdentity_isReduction,
        CZXCompression.czxPlusIdentity_left_mul_right_of_ne,
        CZXCompression.czxPlusIdentity_dim_eq,
        CZXCompression.czxPlusIdentity_evalWord_remainder_eq_zero}
    \leanok
    \uses{def:asymex_czx_plus, thm:asym_multiblock}
    The source of Definition~\ref{def:asymex_czx_plus} carries the block
    triangular form of Theorem~\ref{thm:asym_multiblock} for the four
    slots with $z=3$ zero slots and $9=1+1+2+2+3$. The compression pairs
    are mutually biorthogonal, $W_sV_t=\id$ for $s=t$ and $0$ otherwise,
    they satisfy $W_sB^wV_s=C_s^w$ for every word, and a product of seven
    remainder matrices vanishes.
\end{theorem}

\begin{proof}
    \leanok
    \uses{thm:asymex_explicit_gauge}
    Every letter of $N$ is block diagonal for the splitting
    $\C^3=\C\oplus\C^2$, so every letter of the stacked tensor is block
    diagonal for the induced four-fold splitting of $\C^3\otimes\C^3$.
    Three of those four summands already carry the simple blocks $\delta$,
    $M$ and $M$; the fourth is the stacked tensor of
    Section~\ref{sec:asymex_czx}, whose flag contributes the weight $-1$
    together with the three zero slots. Assembling the corresponding change
    of bond coordinates and checking block triangularity together with the
    matched and unmatched diagonal blocks is a finite verification over
    integer matrices. The dimension count and the nilpotency order are
    clauses (vii) and (vi) of Theorem~\ref{thm:asym_multiblock} with
    $|S|+z=7$.
\end{proof}

\begin{theorem}[The oscillating structure coefficient]
    \label{thm:asymex_czx_plus_trace}
    \lean{CZXCompression.czxPlusIdentity_trace_evalWord}
    \leanok
    \uses{def:asymex_czx_plus}
    For every nonempty word $w$ over the alphabet of index pairs,
    \begin{align}
        \tr(B^w) & =\bigl(1+(-1)^{|w|}\bigr)\tr(\delta^w)+2\tr(M^w).
        \label{eq:asymex_czx_plus_trace}
    \end{align}
\end{theorem}

\begin{proof}
    \leanok
    \uses{thm:asymex_czx_plus_compression, prop:asym_compression_trace}
    Proposition~\ref{prop:asym_compression_trace} applied to the datum of
    Theorem~\ref{thm:asymex_czx_plus_compression} writes $\tr(B^w)$ as the
    sum of the four slot traces; pulling the weight $-1$ out of the second
    slot contributes $(-1)^{|w|}\tr(\delta^w)$ and the two equal slots
    carrying $M$ contribute $2\tr(M^w)$.
\end{proof}

%
\section{Kramers--Wannier duality squared}%
\label{sec:asymex_kw}

The unnormalized Kramers--Wannier duality kernel on a periodic chain of $L$
qubits is
\begin{align}
	\bra{s'}D\ket{s} & =\prod_{j=1}^{L}(-1)^{s'_j(s_j+s_{j+1})},
	\qquad s_{L+1}=s_1,
	\label{eq:asymex_kw_kernel}
\end{align}
the periodic operator of a matrix product operator of bond dimension two.
Squaring the duality gives
\begin{align}
	D^2 & =2^L(1+\eta)\,T,
	\label{eq:asymex_kw_fusion}
\end{align}
where $\eta=\prod_jX_j$ is the global spin flip and $T$ is the one-site
left translation. This is the lattice form of the duality-defect fusion
rule of the Ising model \cite{Aasen2016Topological}, in the normalization in
which the analogous Majorana-chain defect satisfies $D^2=(1+\eta)T$
\cite{Seiberg2024Majorana}. It is a fusion rule of a non-invertible
symmetry whose right-hand side contains a space-time symmetry, so it lies
outside every fusion-category algebra of matrix product operators. Both
$T$ and $\eta T$ are bond-dimension-two matrix product operators whose four
matrices are the four matrix units of $\MN{2}$, hence normal already at
word length one.

\begin{definition}[The duality kernel and the two translation tensors]
	\label{def:asymex_kw}
	\lean{KWExample.kwTensor, KWExample.shiftTensor, KWExample.flipShiftTensor,
		KWExample.kwSquare, KWExample.kwSquareTargets}
	\leanok
	\uses{def:mpo_tensor, def:mpo_to_mps, def:mpdo_operator_product}
	With the outgoing label $s'$ written first, the duality tensor is
	\begin{align}
		W^{00} & =\begin{pmatrix}1&1\\0&0\end{pmatrix},
		       &
		W^{01} & =\begin{pmatrix}0&0\\1&1\end{pmatrix},
		       &
		W^{10} & =\begin{pmatrix}1&-1\\0&0\end{pmatrix},
		       &
		W^{11} & =\begin{pmatrix}0&0\\-1&1\end{pmatrix}.
		\notag
	\end{align}
	The translation tensor is $T^{s's}=E_{s,s'}$ and the flipped
	translation tensor is $(\eta T)^{s's}=E_{s,1-s'}$, where $E_{pq}$ is
	the matrix unit with a single entry $1$ in row $p$ and column $q$. The
	source is the stacked product $B^{s's}=\sum_mW^{s'm}\otimes W^{ms}$ of
	bond dimension four, and the two targets are $2T$ and $2\eta T$.
\end{definition}

\begin{center}
	% Formula: B^{s's}=\sum_mW^{s'm}\otimes W^{ms}, the stacked product
	% of Definition~\ref{def:asymex_kw}.
	% Source: Definition~\ref{def:asymex_kw}.
	% Ink-to-index map: both rows are the duality tensor $W$; the shared
	% vertical pairleg at each site is the summed index $m$; the drawn
	% upper-row label $r$ is the per-site component of the outgoing
	% index $s'$.
	% Contraction set: the vertical pairlegs (index $m$) at every site
	% and the horizontal virtual bonds within each row; the physical
	% legs $(r,s)$, the per-site components of $(s',s)$, remain open.
	% Boundary signature: (virtual-west=0 | virtual-east=0 |
	% physical-up=3 | physical-down=3) on both sides.
	\tenkzkernel
	\begin{tenkz}[rows={op,op}, cols=4, boundary=periodic]
		\tn[skin=mpo, ports={90:physical:$r_1$, 270:physical}]{W} &
		\tn[skin=mpo, ports={90:physical:$r_2$, 270:physical}]{W} &
		\tn[skin=dots, ports={180:virtual, 0:virtual}]{} &
		\tn[skin=mpo, ports={90:physical:$r_N$, 270:physical}]{W} \\
		\tn[skin=mpo, ports={90:physical, 270:physical:$s_1$}]{W} &
		\tn[skin=mpo, ports={90:physical, 270:physical:$s_2$}]{W} &
		\tn[skin=dots, ports={180:virtual, 0:virtual}]{} &
		\tn[skin=mpo, ports={90:physical, 270:physical:$s_N$}]{W}
	\end{tenkz}
	\quad $=$ \quad
	\begin{tenkz}[rows={op}, cols=4, boundary=periodic]
		\tn[skin=mpo, ports={90:physical:$r_1$, 270:physical:$s_1$}]{B} &
		\tn[skin=mpo, ports={90:physical:$r_2$, 270:physical:$s_2$}]{B} &
		\tn[skin=dots, ports={180:virtual, 0:virtual}]{} &
		\tn[skin=mpo, ports={90:physical:$r_N$, 270:physical:$s_N$}]{B}
	\end{tenkz}
\end{center}

\begin{theorem}[Normality of the two translation tensors]
	\label{thm:asymex_kw_normal}
	\lean{KWExample.shiftMPS_isNormal, KWExample.flipShiftMPS_isNormal}
	\leanok
	\uses{def:asymex_kw, def:normal}
	The translation tensor and the flipped translation tensor are normal:
	the four matrices of each are the four matrix units of $\MN{2}$, which
	already span $\MN{2}$ at word length one.
\end{theorem}

\begin{proof}
	\leanok
	Every matrix unit occurs among the four matrices of each tensor, and
	the matrix units span $\MN{2}$.
\end{proof}

\begin{theorem}[Split compression of the squared duality]
	\label{thm:asymex_kw_compression}
	% Principal declaration: KWExample.kwSquare_isReduction.
	\lean{KWExample.kwSquare_compression, KWExample.kwSquare_isReduction,
		KWExample.kwSquare_left_mul_right_of_ne, KWExample.kwSquare_dim_eq,
		KWExample.kwSquare_remainder_eq_zero}
	\leanok
	\uses{def:asymex_kw, thm:asym_multiblock}
	The source of Definition~\ref{def:asymex_kw} carries the block
	triangular form of Theorem~\ref{thm:asym_multiblock} for the two slots
	$2T$ and $2\eta T$ with no zero slots, $4=2+2$. The compression pairs
	are mutually biorthogonal and satisfy $W_sB^wV_s=C_s^w$ for every word,
	and the remainder $R^{s's}=B^{s's}-\sum_sV_sC_s^{s's}W_s$ vanishes
	identically: the extension splits.
\end{theorem}

\begin{proof}
	\leanok
	\uses{thm:asymex_explicit_gauge}
	The change of bond coordinates is
	\begin{align}
		G & =\frac12\begin{pmatrix}
			            1 & 0 & 1  & 0  \\
			            0 & 1 & 0  & -1 \\
			            1 & 0 & -1 & 0  \\
			            0 & 1 & 0  & 1
		            \end{pmatrix},
		\label{eq:asymex_kw_gauge}
	\end{align}
	whose inverse has integer entries. Conjugating each of the four letters
	by $G$ produces the block diagonal matrix with the blocks $2T$ and
	$2\eta T$, which after clearing the denominator is a finite verification
	over integer matrices. Block diagonality gives both the triangularity
	condition and the vanishing of the remainder, and the dimension count is
	clause (vii) of Theorem~\ref{thm:asym_multiblock}.
\end{proof}

\begin{theorem}[Periodic data of the squared duality]
	\label{thm:asymex_kw_trace}
	\lean{KWExample.kwSquare_trace_evalWord}
	\leanok
	\uses{def:asymex_kw}
	For every nonempty word $w$ over the alphabet of index pairs,
	\begin{align}
		\tr(B^w) & =2^{|w|}\bigl(\tr(T^w)+\tr((\eta T)^w)\bigr),
		\label{eq:asymex_kw_trace}
	\end{align}
	which is (\ref{eq:asymex_kw_fusion}) at the level of periodic
	coefficients.
\end{theorem}

\begin{proof}
	\leanok
	\uses{thm:asymex_kw_compression, prop:asym_compression_trace}
	Proposition~\ref{prop:asym_compression_trace} applied to the datum of
	Theorem~\ref{thm:asymex_kw_compression} gives
	$\tr(B^w)=\tr((2T)^w)+\tr((2\eta T)^w)$, and pulling the weight $2$ out
	of each word evaluation contributes the factor $2^{|w|}$.
\end{proof}

\begin{theorem}[The sitewise intertwiners of the split case]
	\label{thm:asymex_kw_intertwiners}
	% Principal declaration: KWExample.kwSquare_mul_right0_eq.
	\lean{KWExample.kwSquareLeft0, KWExample.kwSquareLeft1,
		KWExample.kwSquareRight0, KWExample.kwSquareRight1,
		KWExample.kwSquare_mul_right0_eq, KWExample.kwSquare_mul_right1_eq,
		KWExample.left0_mul_kwSquare_eq, KWExample.left1_mul_kwSquare_eq,
		KWExample.kwSquareLeft0_mul_right0, KWExample.kwSquareLeft1_mul_right1}
	\leanok
	\uses{def:asymex_kw, thm:asymex_kw_compression}
	Write
	\begin{align}
		V_1 & =\begin{pmatrix}
			       1 & 0  \\
			       0 & 1  \\
			       1 & 0  \\
			       0 & -1
		       \end{pmatrix},
		    & V_2
		    & =\begin{pmatrix}
			       1  & 0 \\
			       0  & 1 \\
			       -1 & 0 \\
			       0  & 1
		       \end{pmatrix},
		    & W_1
		    & =\frac12\begin{pmatrix}
			              1 & 0 & 1 & 0  \\
			              0 & 1 & 0 & -1
		              \end{pmatrix},
		    & W_2
		    & =\frac12\begin{pmatrix}
			              1 & 0 & -1 & 0 \\
			              0 & 1 & 0  & 1
		              \end{pmatrix}.
		\notag
	\end{align}
	Then $W_1V_1=W_2V_2=\id$, and for every index pair,
	$B^{s's}V_1=V_1(2T)^{s's}$, $B^{s's}V_2=V_2(2\eta T)^{s's}$,
	$W_1B^{s's}=(2T)^{s's}W_1$ and $W_2B^{s's}=(2\eta T)^{s's}W_2$.
\end{theorem}

\begin{proof}
	\leanok
	\uses{thm:asymex_explicit_gauge}
	These four matrices are the rows of $G$ and the columns of $G^{-1}$ at
	the coordinates of the two slots, by
	Theorem~\ref{thm:asymex_explicit_gauge}. Each of the four intertwining
	identities is an equality between products of integer matrices after
	clearing the denominator, and the two normalizations are clause (iv) of
	Theorem~\ref{thm:asym_multiblock}.
\end{proof}

\begin{center}
	% Formula: the four sitewise identities of
	% Theorem~\ref{thm:asymex_kw_intertwiners}: B^{s's}V_1=V_1(2T)^{s's},
	% B^{s's}V_2=V_2(2\eta T)^{s's}, W_1B^{s's}=(2T)^{s's}W_1, and
	% W_2B^{s's}=(2\eta T)^{s's}W_2.
	% Source: Theorem~\ref{thm:asymex_kw_intertwiners}.
	% Ink-to-index map: the ring is the compression map $V_1$, $W_1$,
	% $V_2$ or $W_2$; the square node is the source $B$ or its target
	% $2T$ or $2\eta T$.
	% Contraction set: the single virtual bond joining the ring to the
	% square node in each pair; the physical pair leg remains open.
	% Boundary signature: (virtual-west=1 | virtual-east=1 |
	% physical-up=1 | physical-down=0) for each pictured equality.
	\tenkzkernel
	\begin{tenkz}[rows={wire}, cols=2, west=open, east=open]
		\tn[label pos=s, ports={90:physical:$s's$}]{B} & \tn[skin=ring]{V_1}
	\end{tenkz}
	\quad $=$ \quad
	\begin{tenkz}[rows={wire}, cols=2, west=open, east=open]
		\tn[skin=ring]{V_1} & \tn[label pos=s, ports={90:physical:$s's$}]{2T}
	\end{tenkz}
	\\[2ex]
	\begin{tenkz}[rows={wire}, cols=2, west=open, east=open]
		\tn[label pos=s, ports={90:physical:$s's$}]{B} & \tn[skin=ring]{V_2}
	\end{tenkz}
	\quad $=$ \quad
	\begin{tenkz}[rows={wire}, cols=2, west=open, east=open]
		\tn[skin=ring]{V_2} &
		\tn[label pos=s, ports={90:physical:$s's$}]{2\eta T}
	\end{tenkz}
	\\[2ex]
	\begin{tenkz}[rows={wire}, cols=2, west=open, east=open]
		\tn[skin=ring]{W_1} & \tn[label pos=s, ports={90:physical:$s's$}]{B}
	\end{tenkz}
	\quad $=$ \quad
	\begin{tenkz}[rows={wire}, cols=2, west=open, east=open]
		\tn[label pos=s, ports={90:physical:$s's$}]{2T} & \tn[skin=ring]{W_1}
	\end{tenkz}
	\\[2ex]
	\begin{tenkz}[rows={wire}, cols=2, west=open, east=open]
		\tn[skin=ring]{W_2} & \tn[label pos=s, ports={90:physical:$s's$}]{B}
	\end{tenkz}
	\quad $=$ \quad
	\begin{tenkz}[rows={wire}, cols=2, west=open, east=open]
		\tn[label pos=s, ports={90:physical:$s's$}]{2\eta T} &
		\tn[skin=ring]{W_2}
	\end{tenkz}
\end{center}

\section{Kramers--Wannier duality acting on ground states}%
\label{sec:asymex_kw_action}

The same duality kernel may be applied to a state rather than to itself.
Both parts below read as symmetry breaking of the non-invertible
Kramers--Wannier symmetry: the duality maps the paramagnetic product state
to the pair of ferromagnetic sectors, and it maps both ferromagnetic
sectors to the paramagnet, merged with multiplicity two. When the acted
state is written as a multiple of a normalized state, a factor $2^{L/2}$
appears; that factor is a norm of the unnormalized periodic vector and not
a weight of the target, and the weighted tensor seen by the theorem is the
same under either reading.

\begin{definition}[The two action tensors]
	\label{def:asymex_kw_action}
	\lean{KWExample.plusTensor, KWExample.ghz0, KWExample.ghz1,
		KWExample.kwPlus, KWExample.plusTarget, KWExample.ghz,
		KWExample.kwGHZ, KWExample.kwGHZTarget}
	\leanok
	\uses{def:asymex_kw, def:mps_tensor}
	The product-state tensor is the bond-dimension-one tensor with both
	matrices equal to $(1)$, generating $\sum_s\ket{s}$ at every site; the
	two ferromagnetic sectors are the bond-dimension-one tensors
	$C_0=\bigl((1),(0)\bigr)$ and $C_1=\bigl((0),(1)\bigr)$, and their direct
	sum is the Greenberger--Horne--Zeilinger tensor $\diag(1,0)$,
	$\diag(0,1)$.
	The two action tensors are
	$(W\cdot A)^{s'}=\sum_sW^{s's}\otimes A^s$ applied to the product-state
	tensor and to the Greenberger--Horne--Zeilinger tensor; explicitly
	\begin{align}
		(W\cdot A)^0
		 & =\begin{pmatrix}1&1\\1&1\end{pmatrix},
		 &
		(W\cdot A)^1
		 & =\begin{pmatrix}1&-1\\-1&1\end{pmatrix}
		\notag
	\end{align}
	in the first case, with targets the two sectors carried with weight $2$,
	and
	\begin{align}
		(W\cdot A)^0
		 & =\begin{pmatrix}
			    1 & 0 & 1 & 0 \\
			    0 & 0 & 0 & 0 \\
			    0 & 0 & 0 & 0 \\
			    0 & 1 & 0 & 1
		    \end{pmatrix},
		 &
		(W\cdot A)^1
		 & =\begin{pmatrix}
			    1 & 0  & -1 & 0 \\
			    0 & 0  & 0  & 0 \\
			    0 & 0  & 0  & 0 \\
			    0 & -1 & 0  & 1
		    \end{pmatrix}
		\notag
	\end{align}
	in the second, with two unweighted copies of the product-state tensor as
	targets.
\end{definition}

\begin{center}
	% Formula: (W\cdot A)^{s'}=\sum_sW^{s's}\otimes A^s, of
	% Definition~\ref{def:asymex_kw_action}, evaluated at the
	% product-state tensor $A_+$ (left) and at the Greenberger--Horne--
	% Zeilinger tensor $A_{\mathrm{GHZ}}$ (right).
	% Source: Definition~\ref{def:asymex_kw_action}.
	% Ink-to-index map: the upper row is the duality tensor $W$, whose
	% drawn physical legs $r$ are the per-site components of the
	% outgoing index $s'$; the lower row is the acted-on state tensor;
	% the shared vertical pairleg at each site is the summed index $s$.
	% Contraction set: the vertical pairlegs (index $s$) at every site
	% and the horizontal virtual bonds within each row, both closed by
	% the periodic trace; only the upper physical legs of $W$ remain
	% open.
	% Boundary signature: (virtual-west=0 | virtual-east=0 |
	% physical-up=3 | physical-down=0) for both pictures.
	\tenkzkernel
	\begin{tenkz}[rows={op,ket}, cols=4, boundary=periodic]
		\tn[skin=mpo, ports={90:physical:$r_1$}]{W} &
		\tn[skin=mpo, ports={90:physical:$r_2$}]{W} &
		\tn[skin=dots, ports={180:virtual, 0:virtual}]{} &
		\tn[skin=mpo, ports={90:physical:$r_N$}]{W} \\
		\tn[label pos=s, ports={180:virtual, 0:virtual}]{A_+} &
		\tn[label pos=s, ports={180:virtual, 0:virtual}]{A_+} &
		\tn[skin=dots, ports={180:virtual, 0:virtual}]{} &
		\tn[label pos=s, ports={180:virtual, 0:virtual}]{A_+}
	\end{tenkz}
	\qquad
	\begin{tenkz}[rows={op,ket}, cols=4, boundary=periodic]
		\tn[skin=mpo, ports={90:physical:$r_1$}]{W} &
		\tn[skin=mpo, ports={90:physical:$r_2$}]{W} &
		\tn[skin=dots, ports={180:virtual, 0:virtual}]{} &
		\tn[skin=mpo, ports={90:physical:$r_N$}]{W} \\
		\tn[label pos=s, ports={180:virtual, 0:virtual}]{A_{\mathrm{GHZ}}} &
		\tn[label pos=s, ports={180:virtual, 0:virtual}]{A_{\mathrm{GHZ}}} &
		\tn[skin=dots, ports={180:virtual, 0:virtual}]{} &
		\tn[label pos=s, ports={180:virtual, 0:virtual}]{A_{\mathrm{GHZ}}}
	\end{tenkz}
\end{center}

\begin{theorem}[The duality on the product state splits]
	\label{thm:asymex_kw_plus}
	% Principal declaration: KWExample.plus_isReduction.
	\lean{KWExample.plusCompression, KWExample.plus_isReduction,
		KWExample.plus_dim_eq, KWExample.plus_remainder_eq_zero,
		KWExample.kwPlus_trace_evalWord_eq_sum}
	\leanok
	\uses{def:asymex_kw_action, thm:asym_multiblock}
	The first action tensor of Definition~\ref{def:asymex_kw_action} carries
	the block triangular form of Theorem~\ref{thm:asym_multiblock} for the
	two sectors carried with weight $2$, with no zero slots and $2=1+1$;
	the compression pair of each sector satisfies $W_sV_s=\id$ and
	$W_sB^wV_s=C_s^w$ for every word; the remainder vanishes identically;
	and for every nonempty word $w$,
	$\tr(B^w)=\tr((2C_0)^w)+\tr((2C_1)^w)$.
\end{theorem}

\begin{proof}
	\leanok
	\uses{thm:asymex_explicit_gauge, prop:asym_compression_trace}
	Conjugating the two letters by the matrix
	$\begin{pmatrix}1&1\\1&-1\end{pmatrix}$, whose inverse is half of
	itself, produces $\diag(2,0)$ and $\diag(0,2)$, a finite verification
	over integer matrices. Diagonality gives the triangularity condition,
	the identification of the matched blocks and the vanishing of the
	remainder; the trace identity is
	Proposition~\ref{prop:asym_compression_trace}.
\end{proof}

\begin{theorem}[The duality on the ferromagnetic state does not split]
	\label{thm:asymex_kw_ghz}
	% Principal declaration: KWExample.kwGHZ_isReduction.
	\lean{KWExample.kwGHZCompression, KWExample.kwGHZ_isReduction,
		KWExample.kwGHZ_dim_eq, KWExample.kwGHZ_evalWord_remainder_eq_zero,
		KWExample.kwGHZRemainder_eq, KWExample.kwGHZ_remainder_mul_remainder,
		KWExample.kwGHZ_trace_evalWord_eq_two}
	\leanok
	\uses{def:asymex_kw_action, thm:asym_multiblock}
	The second action tensor of Definition~\ref{def:asymex_kw_action}
	carries the block triangular form of Theorem~\ref{thm:asym_multiblock}
	for two unweighted copies of the product-state tensor with $z=2$ zero
	slots and $4=1+1+2$; the compression pair of each copy satisfies
	$W_sV_s=\id$ and $W_sB^wV_s=C_s^w$ for every word;
	$\tr(B^w)=2$ for every nonempty word $w$; a product of four remainder
	matrices vanishes; and in fact
	\begin{align}
		R^0 & =\begin{pmatrix}
			       0 & 0 & 1 & 0 \\
			       0 & 0 & 0 & 0 \\
			       0 & 0 & 0 & 0 \\
			       0 & 1 & 0 & 0
		       \end{pmatrix},
		    & R^1
		    & =\begin{pmatrix}
			       0 & 0  & -1 & 0 \\
			       0 & 0  & 0  & 0 \\
			       0 & 0  & 0  & 0 \\
			       0 & -1 & 0  & 0
		       \end{pmatrix},
		\notag
	\end{align}
	and $R^iR^j=0$ for all letters $i$ and $j$.
\end{theorem}

\begin{proof}
	\leanok
	\uses{thm:asymex_explicit_gauge, prop:asym_compression_trace}
	The gauge is the coordinate permutation placing the two copies at the
	first and fourth bond coordinates and the two zero slots at the second
	and third; in these coordinates the two letters are block upper
	triangular with the scalar $1$ on both matched diagonal positions, a
	finite verification over integer matrices. Every word evaluation of the
	product-state tensor is $1$, so the trace identity of
	Proposition~\ref{prop:asym_compression_trace} reads $\tr(B^w)=2$. The
	remainder is the conjugated letter off the diagonal of the block
	decomposition, which gives the two displayed matrices, and their
	pairwise products vanish by direct multiplication; the generic bound of
	clause (vi) of Theorem~\ref{thm:asym_multiblock} is $|S|+z=4$.
\end{proof}

\begin{theorem}[One-sided sitewise intertwiners for the merged target]
	\label{thm:asymex_kw_ghz_intertwiner}
	\lean{KWExample.kwGHZ_left_intertwiner_eq_zero,
		KWExample.kwGHZ_right_intertwiner,
		KWExample.kwGHZ_right_intertwiner_ne_zero}
	\leanok
	\uses{def:asymex_kw_action}
	If $u\in\C^4$ satisfies $uB^i=u$ for both letters $i$, then $u=0$: no
	nonzero sitewise left intertwiner exists for the product-state target.
	A nonzero sitewise right intertwiner does exist, namely the first
	coordinate vector, which satisfies $B^ie_0=e_0$ for both letters.
\end{theorem}

\begin{proof}
	\leanok
	Both matrices of the product-state target are the scalar $1$, so the
	two conditions read $uB^i=u$ and $B^iv=v$. Comparing the second and
	third coordinates of the two equations $uB^0=u$ and $uB^1=u$ gives in
	turn $u_3=u_1=u_0=u_2=0$. For the right-hand side, the first column of
	each of the two letters is the first coordinate vector.
\end{proof}

\section{Fibonacci fusion}%
\label{sec:asymex_fibonacci}

The simplest non-invertible fusion of matrix product operators is the
Fibonacci string-net algebra of \cite[Appendix~D.1]{Bultinck2017Anyons},
whose two operator families obey
\begin{align}
	O_\tau O_\tau & =O_1+O_\tau,
	\label{eq:asymex_fib_fusion}
\end{align}
the realization of $\tau\otimes\tau=1\oplus\tau$. The unit of the algebra
is not the identity operator: $O_1$ is the projector onto the admissible
configurations of the chain, those with no two neighbouring trivial labels,
and it is a matrix product operator of bond dimension two. The $\tau$
family has bond dimension three, so the stacked tensor of the square has
bond dimension nine: two directions carry the trivial block, three carry
the $\tau$ block, and the remaining four are annihilated by every letter,
being the inadmissible configurations killed by the vertex weights.

The entries of these tensors are not rational, so the exact arithmetic
takes place in a quartic ring.

\begin{definition}[The ring of golden integers]
	\label{def:asymex_golden_ring}
	\lean{GoldenInt, GoldenInt.sigma, goldenPoly, goldenSigmaReal,
		goldenToComplex, MPSTensor.complexOfGolden, MPSTensor.evalWordGolden,
		MPSTensor.mulGoldenTensor}
	\leanok
	Let $\sigma=\varphi^{-1/2}$ be the inverse square root of the golden
	ratio, the positive real root of $X^4+X^2-1$, so that
	$\sigma=\sqrt{(\sqrt5-1)/2}$. The ring of \emph{golden integers} is
	$\Z[\sigma]$, presented by the four integer coordinates of
	$c_0+c_1\sigma+c_2\sigma^2+c_3\sigma^3$ with the multiplication obtained
	from $\sigma^4=1-\sigma^2$. It embeds into $\C$ through the quotient
	$\mathbb{Q}[X]/(X^4+X^2-1)$ evaluated at that real root; a matrix over
	$\Z[\sigma]$ becomes a complex matrix by applying that embedding
	entrywise, and word evaluation and the bond-space product of tensors
	over $\Z[\sigma]$ are defined by the same formulas as over $\C$.
\end{definition}

\begin{theorem}[Exact arithmetic in the ring of golden integers]
	\label{thm:asymex_golden_ring}
	% Principal declaration: goldenMulRule.
	\lean{GoldenInt.sigma_quartic, goldenMulRule, goldenSigmaReal_quartic,
		MPSTensor.complexOfGolden_mul, MPSTensor.evalWord_complexOfGolden,
		MPSTensor.mulTensor_complexOfGolden}
	\leanok
	\uses{def:asymex_golden_ring}
	The generator satisfies $\sigma^4+\sigma^2-1=0$ in $\Z[\sigma]$, and so
	does the real number $\sqrt{(\sqrt5-1)/2}$. In any commutative ring
	carrying a root $s$ of $X^4+X^2-1$, the product of two integer
	combinations of $1,s,s^2,s^3$ is the combination prescribed by the
	multiplication of $\Z[\sigma]$. The entrywise embedding of matrices
	over $\Z[\sigma]$ into the complex matrices is multiplicative, and it
	commutes with word evaluation and with the bond-space product of
	tensors.
\end{theorem}

\begin{proof}
	\leanok
	Squaring $\sigma^2=(\sqrt5-1)/2$ gives the quartic for the real root.
	The multiplication rule is the reduction of a product of two
	polynomials of degree at most three modulo $X^4+X^2-1$, which is a
	polynomial identity in the eight coefficients. Multiplicativity of the
	entrywise embedding is the fact that it is induced by a ring
	homomorphism, and the two compatibility statements follow by induction
	on the word and by expanding the bond-space product entrywise.
\end{proof}

\begin{definition}[The two Fibonacci blocks and the stacked square]
	\label{def:asymex_fib}
	\lean{FibonacciCompression.fibTau, FibonacciCompression.fibOne,
		FibonacciCompression.fibTauMPS, FibonacciCompression.fibOneMPS,
		FibonacciCompression.fibStack, FibonacciCompression.fibTargets}
	\leanok
	\uses{def:asymex_golden_ring, def:mpo_tensor, def:mpdo_operator_product}
	Label the two objects of the category by $0$ for the trivial label and
	$1$ for $\tau$, and read a letter as the pair $(x',x)$ of the outgoing
	and incoming labels. Write
	\begin{align}
		F & =\begin{pmatrix}
			     0 & \varphi^{-1}   & \varphi^{-1/2} \\
			     1 & 0              & 1              \\
			     1 & \varphi^{-1/2} & -\varphi^{-1}
		     \end{pmatrix}
		\label{eq:asymex_fib_vertex}
	\end{align}
	for the reduced vertex weight. The $\tau$ tensor $B_\tau$ of bond
	dimension three has, at the letters $(0,1)$, $(1,0)$ and $(1,1)$, the
	matrix whose first, second and third row respectively is the
	corresponding row of $F$ and whose other rows vanish; at the letter
	$(0,0)$ it vanishes. The trivial tensor $B_1$ of bond dimension two
	carries the admissibility matrix
	$\begin{pmatrix}0&1\\1&1\end{pmatrix}$ on its two diagonal letters and
	vanishes on the other two. The source is the stacked product
	$B^{x'x}=\sum_m(B_\tau)^{x'm}\otimes(B_\tau)^{mx}$ of bond dimension
	nine, and the two targets are $B_1$ and $B_\tau$, each with
	multiplicity one and weight one.
\end{definition}

\begin{center}
	% Formula: B^{x'x}=\sum_m(B_\tau)^{x'm}\otimes(B_\tau)^{mx}, the
	% stacked product of Definition~\ref{def:asymex_fib}, realizing
	% $\tau\otimes\tau$; Theorem~\ref{thm:asymex_fib_fusion} identifies
	% its periodic trace with that of $B_1\oplus B_\tau$, the
	% realization $1\oplus\tau$.
	% Source: Definition~\ref{def:asymex_fib}.
	% Ink-to-index map: both rows are the $\tau$ tensor $B_\tau$; the
	% shared vertical pairleg at each site is the summed intermediate
	% label $m$; the drawn upper-row label $y$ is the per-site
	% component of the outgoing index $x'$.
	% Contraction set: the vertical pairlegs (index $m$) at every site
	% and the horizontal virtual bonds within each row; the physical
	% legs $(y,x)$, the per-site components of $(x',x)$, remain open.
	% Boundary signature: (virtual-west=0 | virtual-east=0 |
	% physical-up=3 | physical-down=3) on both sides.
	\tenkzkernel
	\begin{tenkz}[rows={op,op}, cols=4, boundary=periodic]
		\tn[skin=mpo, ports={90:physical:$y_1$, 270:physical}]{B_\tau} &
		\tn[skin=mpo, ports={90:physical:$y_2$, 270:physical}]{B_\tau} &
		\tn[skin=dots, ports={180:virtual, 0:virtual}]{} &
		\tn[skin=mpo, ports={90:physical:$y_N$, 270:physical}]{B_\tau} \\
		\tn[skin=mpo, ports={90:physical, 270:physical:$x_1$}]{B_\tau} &
		\tn[skin=mpo, ports={90:physical, 270:physical:$x_2$}]{B_\tau} &
		\tn[skin=dots, ports={180:virtual, 0:virtual}]{} &
		\tn[skin=mpo, ports={90:physical, 270:physical:$x_N$}]{B_\tau}
	\end{tenkz}
	\quad $=$ \quad
	\begin{tenkz}[rows={op}, cols=4, boundary=periodic]
		\tn[skin=mpo, ports={90:physical:$y_1$, 270:physical:$x_1$}]{B} &
		\tn[skin=mpo, ports={90:physical:$y_2$, 270:physical:$x_2$}]{B} &
		\tn[skin=dots, ports={180:virtual, 0:virtual}]{} &
		\tn[skin=mpo, ports={90:physical:$y_N$, 270:physical:$x_N$}]{B}
	\end{tenkz}
\end{center}

\begin{theorem}[The one-site matrices of the $\tau$ block]
	\label{thm:asymex_fib_entries}
	\lean{FibonacciCompression.fibTauMPS_apply_one,
		FibonacciCompression.fibTauMPS_apply_three}
	\leanok
	\uses{def:asymex_fib}
	The matrix of $B_\tau$ at the letter $(0,1)$ has the first row of
	(\ref{eq:asymex_fib_vertex}) as its first row and zero elsewhere, and
	its matrix at the letter $(1,1)$ has the third row of
	(\ref{eq:asymex_fib_vertex}) as its third row and zero elsewhere, with
	$\varphi^{-1/2}=\sqrt{(\sqrt5-1)/2}$.
\end{theorem}

\begin{proof}
	\leanok
	\uses{thm:asymex_golden_ring}
	Both are the entrywise embedding of the corresponding matrix over
	$\Z[\sigma]$, evaluated coordinate by coordinate.
\end{proof}

\begin{theorem}[Normality of the two Fibonacci blocks]
	\label{thm:asymex_fib_normal}
	\lean{FibonacciCompression.fibOneMPS_isNormal,
		FibonacciCompression.fibTauMPS_isNormal}
	\leanok
	\uses{def:asymex_fib, def:normal}
	Both blocks are normal, their length-three words spanning the full
	matrix algebra of their bond space.
\end{theorem}

\begin{proof}
	\leanok
	\uses{thm:asymex_golden_ring}
	For the $\tau$ block the word $(i+1,3,k+1)$, with $i$ and $k$ in
	$\{0,1,2\}$, evaluates to a multiple of the matrix built from the $i$-th
	coordinate vector and the $k$-th row of (\ref{eq:asymex_fib_vertex});
	those rows are a basis, so each of the nine matrix units is an explicit
	combination over $\Z[\sigma]$ of the nine such words. For the trivial block four length-three words suffice. Both
	families of coefficients are verified in $\Z[\sigma]$ and transported to
	the complex matrices by Theorem~\ref{thm:asymex_golden_ring}.
\end{proof}

\begin{theorem}[Split compression of the Fibonacci square]
	\label{thm:asymex_fib_compression}
	% Principal declaration: FibonacciCompression.fibonacci_isReduction.
	\lean{FibonacciCompression.fibonacci_compression,
		FibonacciCompression.fibonacci_isReduction,
		FibonacciCompression.fibonacci_left_mul_right_of_ne,
		FibonacciCompression.fibonacci_dim_eq,
		FibonacciCompression.fibonacci_evalWord_remainder_eq_zero,
		FibonacciCompression.fibonacci_remainder_eq_zero}
	\leanok
	\uses{def:asymex_fib, thm:asym_multiblock}
	The source of Definition~\ref{def:asymex_fib} carries the block
	triangular form of Theorem~\ref{thm:asym_multiblock} for the trivial
	block and the $\tau$ block with $z=4$ zero slots and $9=2+3+4$. The
	compression pairs are mutually biorthogonal and satisfy
	$W_sB^wV_s=C_s^w$ for every word; a product of six remainder matrices
	vanishes; and in fact the remainder vanishes identically, so the
	extension splits.
\end{theorem}

\begin{proof}
	\leanok
	\uses{thm:asymex_explicit_gauge, thm:asymex_golden_ring}
	The change of bond coordinates has as its columns the two fusion
	tensors of the blocks followed by a basis of the four-dimensional
	subspace annihilated by every letter. In these coordinates every letter
	of the stacked tensor is block diagonal with the trivial block, the
	$\tau$ block and four one-dimensional zero blocks on the diagonal, which
	is a finite verification over matrices with entries in $\Z[\sigma]$.
	Block diagonality gives the triangularity condition, the identification
	of the diagonal blocks and the vanishing of the remainder; the generic
	nilpotency bound of clause (vi) of Theorem~\ref{thm:asym_multiblock} is
	$|S|+z=6$, and the dimension count is clause (vii).
\end{proof}

\begin{theorem}[The fusion rule]
	\label{thm:asymex_fib_fusion}
	% Principal declaration: FibonacciCompression.fibonacci_fusion_rule.
	\lean{FibonacciCompression.fibStack_trace_evalWord,
		FibonacciCompression.fibonacci_fusion_rule,
		FibonacciCompression.fibStack_mpv}
	\leanok
	\uses{def:asymex_fib, def:mpo_family}
	For every nonempty word $w$ over the alphabet of label pairs,
	$\tr(B^w)=\tr(B_1^w)+\tr(B_\tau^w)$; equivalently, at every positive
	system size $L$ the periodic operators satisfy
	(\ref{eq:asymex_fib_fusion}), and the periodic vector of the stacked
	tensor is the sum of the periodic vectors of the two blocks at every
	positive size.
\end{theorem}

\begin{proof}
	\leanok
	\uses{thm:asymex_fib_compression, prop:asym_compression_trace}
	The word-trace identity is
	Proposition~\ref{prop:asym_compression_trace} applied to the datum of
	Theorem~\ref{thm:asymex_fib_compression}. Each matrix element of a
	periodic operator at size $L$ is the word trace at the word of length
	$L$ formed by the corresponding pairs of labels, and the periodic
	operator of the stacked tensor is the product of the two periodic
	operators, so the identity of periodic operators follows entry by entry;
	the statement for periodic vectors is the same computation.
\end{proof}

\begin{theorem}[The fusion tensors]
	\label{thm:asymex_fib_fusion_tensors}
	% Principal declaration: FibonacciCompression.fibStack_mul_fusionRight.
	\lean{FibonacciCompression.fibFusionRight, FibonacciCompression.fibFusionLeft,
		FibonacciCompression.fibonacci_right_eq, FibonacciCompression.fibonacci_left_eq,
		FibonacciCompression.fibStack_mul_fusionRight,
		FibonacciCompression.fusionLeft_mul_fibStack}
	\leanok
	\uses{def:asymex_fib, thm:asymex_fib_compression}
	The compression maps of the two slots are the fusion tensors of the
	Fibonacci algebra: $V_s$ is the column block of the adapted basis
	carrying the slot $s$, $W_s$ is the matching row block of its inverse,
	and for every letter and every slot,
	$B^{x'x}V_s=V_sC_s^{x'x}$ and $W_sB^{x'x}=C_s^{x'x}W_s$.
\end{theorem}

\begin{proof}
	\leanok
	\uses{thm:asymex_explicit_gauge}
	The identification of the two maps with the row and column blocks of
	the adapted basis is Theorem~\ref{thm:asymex_explicit_gauge}, and each
	of the two intertwining identities is an equality between products of
	matrices over $\Z[\sigma]$, transported to the complex matrices by
	Theorem~\ref{thm:asymex_golden_ring}.
\end{proof}

\begin{center}
	% Formula: B^{x'x}V_s=V_sC_s^{x'x} and W_sB^{x'x}=C_s^{x'x}W_s of
	% Theorem~\ref{thm:asymex_fib_fusion_tensors}, for $s\in\{1,\tau\}$.
	% Source: Theorem~\ref{thm:asymex_fib_fusion_tensors}.
	% Ink-to-index map: the ring is the fusion tensor $V_s$ or $W_s$;
	% the square node is the stacked source $B$ or the target block
	% $C_s\in\{B_1,B_\tau\}$.
	% Contraction set: the single virtual bond joining the ring to the
	% square node in each pair; the physical pair leg $(x',x)$ remains
	% open.
	% Boundary signature: (virtual-west=1 | virtual-east=1 |
	% physical-up=1 | physical-down=0) for each pictured equality.
	\tenkzkernel
	\begin{tenkz}[rows={wire}, cols=2, west=open, east=open]
		\tn[label pos=s, ports={90:physical:$x'x$}]{B} & \tn[skin=ring]{V_s}
	\end{tenkz}
	\quad $=$ \quad
	\begin{tenkz}[rows={wire}, cols=2, west=open, east=open]
		\tn[skin=ring]{V_s} & \tn[label pos=s, ports={90:physical:$x'x$}]{C_s}
	\end{tenkz}
	\\[2ex]
	\begin{tenkz}[rows={wire}, cols=2, west=open, east=open]
		\tn[skin=ring]{W_s} & \tn[label pos=s, ports={90:physical:$x'x$}]{B}
	\end{tenkz}
	\quad $=$ \quad
	\begin{tenkz}[rows={wire}, cols=2, west=open, east=open]
		\tn[label pos=s, ports={90:physical:$x'x$}]{C_s} & \tn[skin=ring]{W_s}
	\end{tenkz}
\end{center}

\section{Parity-graded sectors}%
\label{sec:asymex_parity}

A bond-two normal tensor paired with its own copy of weight $-1$ has the
length-$L$ coefficient $1+(-1)^L$, the power sum of the weight multiset
$\{+1,-1\}$; in the Kitaev chain this grading is the fermion-parity sign
picked up under parity-twisted boundary conditions. The point of the
example is that the source is presented in a mixing integer basis in which
neither of its two matrices has a zero row or a zero column, so the
invariant flag is not visible on the nose, and that the generic nilpotency
bound of the theorem is attained exactly.

\begin{definition}[The graded pair and the five-dimensional source]
	\label{def:asymex_parity}
	\lean{ParityGraded.parA, ParityGraded.parTargets, ParityGraded.parB}
	\leanok
	\uses{def:mps_tensor}
	The block is the bond-two tensor
	\begin{align}
		A^0 & =\begin{pmatrix}1&0\\0&0\end{pmatrix},
		    &
		A^1 & =\begin{pmatrix}0&1\\1&0\end{pmatrix},
		\notag
	\end{align}
	and the two targets are $A$ and $-A$. The source is the
	bond-dimension-five tensor
	\begin{align}
		B^0
		 & =\begin{pmatrix}
			    -1 & 1 & -1 & 1  & 3 \\
			    2  & 2 & -2 & -1 & 0 \\
			    1  & 2 & -2 & 0  & 2 \\
			    1  & 3 & -3 & -1 & 2 \\
			    -1 & 0 & 0  & 1  & 2
		    \end{pmatrix},
		 &
		B^1
		 & =\begin{pmatrix}
			    3  & 1 & 0  & -2 & -3 \\
			    -3 & 1 & -1 & 2  & 5  \\
			    -1 & 2 & -1 & 1  & 3  \\
			    0  & 1 & -1 & -1 & 1  \\
			    2  & 1 & 0  & -1 & -2
		    \end{pmatrix}.
		\notag
	\end{align}
\end{definition}

\begin{center}
	% Formula: the periodic trace tr(B^w)=(1+(-1)^{|w|})tr(A^w) of
	% Theorem~\ref{thm:asymex_parity_trace}; the source is not presented
	% as a literal stack of $A$ and $-A$, so the ring carries $B$ with
	% its coefficient, not a two-row stack.
	% Source: Definition~\ref{def:asymex_parity}.
	% Ink-to-index map: each node is the bond-dimension-five source
	% tensor $B$; the upper legs carry the letters of the word.
	% Contraction set: the horizontal virtual legs, closed around the
	% ring by the trace; no leg is left open.
	% Boundary signature: (virtual-west=0 | virtual-east=0 |
	% physical-up=3 | physical-down=0) for the finite schematic shown.
	\tenkzkernel
	\begin{tenkz}[cols=4, boundary=periodic, physical=up]
		\tn[label pos=135, ports={90:physical:$i_1$}]{B} &
		\tn[label pos=135]{B} & \tn[skin=dots]{} &
		\tn[label pos=135, ports={90:physical:$i_N$}]{B}
	\end{tenkz}
	\qquad $\tr(B^w)$, computed in
	Theorem~\ref{thm:asymex_parity_trace}
\end{center}

\begin{theorem}[Normality of the graded block]
	\label{thm:asymex_parity_normal}
	\lean{ParityGraded.parA_isNormal}
	\leanok
	\uses{def:asymex_parity, def:normal}
	The block $A$ is normal: its length-two words span $\MN{2}$.
\end{theorem}

\begin{proof}
	\leanok
	Three of the four matrix units are directly the products $A^0A^0$,
	$A^0A^1$ and $A^1A^0$, and the fourth is $A^1A^1-A^0A^0$; all four
	identities are finite verifications over integer matrices.
\end{proof}

\begin{theorem}[Compression of the graded pair]
	\label{thm:asymex_parity_compression}
	% Principal declaration: ParityGraded.parityGraded_isReduction.
	\lean{ParityGraded.parityGraded_compression,
		ParityGraded.parityGraded_isReduction,
		ParityGraded.parityGraded_left_mul_right_of_ne,
		ParityGraded.parityGraded_dim_eq,
		ParityGraded.parityGraded_evalWord_remainder_eq_zero,
		ParityGraded.parityGraded_remainder_sq_ne_zero}
	\leanok
	\uses{def:asymex_parity, thm:asym_multiblock}
	The source of Definition~\ref{def:asymex_parity} carries the block
	triangular form of Theorem~\ref{thm:asym_multiblock} for the two slots
	$A$ and $-A$ with $z=1$ zero slot and $5=2+2+1$. The compression pairs
	are mutually biorthogonal and satisfy $W_sB^wV_s=C_s^w$ for every word,
	a product of three remainder matrices vanishes, and the square of the
	remainder at the letter $0$ is nonzero, so the bound $|S|+z=3$ of
	clause (vi) is attained.
\end{theorem}

\begin{proof}
	\leanok
	\uses{thm:asymex_explicit_gauge}
	Conjugating the two letters by an explicit integer matrix with integer
	inverse produces block upper triangular matrices carrying $A$ on the
	first diagonal block, $-A$ on the second, and zero on the third, a
	finite verification over integer matrices. The compression pairs are
	then the rows and columns of that matrix and its inverse at the
	coordinates of the two slots, by
	Theorem~\ref{thm:asymex_explicit_gauge}, and computing the remainder
	from them exhibits a nonzero entry in its square.
\end{proof}

\begin{theorem}[The graded structure constant]
	\label{thm:asymex_parity_trace}
	\lean{ParityGraded.parityGraded_trace_evalWord}
	\leanok
	\uses{def:asymex_parity}
	For every nonempty word $w$,
	$\tr(B^w)=\bigl(1+(-1)^{|w|}\bigr)\tr(A^w)$.
\end{theorem}

\begin{proof}
	\leanok
	\uses{thm:asymex_parity_compression, prop:asym_compression_trace}
	Proposition~\ref{prop:asym_compression_trace} applied to the datum of
	Theorem~\ref{thm:asymex_parity_compression} gives
	$\tr(B^w)=\tr(A^w)+\tr((-A)^w)$, and pulling the weight $-1$ out of the
	second word evaluation contributes $(-1)^{|w|}\tr(A^w)$.
\end{proof}

\begin{theorem}[The weight $-1$ block has no sitewise intertwiner]
	\label{thm:asymex_parity_no_intertwiner}
	\lean{ParityGraded.parNeg_right_intertwiner_eq_zero,
		ParityGraded.parNeg_left_intertwiner_eq_zero}
	\leanok
	\uses{def:asymex_parity}
	If $V\in\C^{5\times2}$ satisfies $B^iV=V(-A^i)$ for both letters, then
	$V=0$; if $W\in\C^{2\times5}$ satisfies $WB^i=(-A^i)W$ for both
	letters, then $W=0$.
\end{theorem}

\begin{proof}
	\leanok
	Each of the two conditions is a linear system in the ten unknown
	entries. Ten of its coordinate equations, combined with explicit
	rational coefficients, determine the five entries of the first column or
	row, and each entry of the remaining one is then read off a further
	equation; both systems have only the zero solution.
\end{proof}

%
\section{Renormalization fixed points: a shared adapted basis}%
\label{sec:asymex_rfp_basis}

The two examples considered below use stacked products of matrix product
operator tensors whose bond spaces are tensor squares. Reordering the
stacked bond coordinates separates a four-dimensional inner bond from
the outer bond space. The inner bond consists of the two half-bonds
joined by summing over the intermediate physical index. Both examples
use the same adapted basis: two Bell vectors and the two coordinate
vectors orthogonal to them.

The length dependence comes from powers of the weights attached to
rank-one idempotents. In the notation of the RFP algebra, the structure
coefficients are
$c_{\alpha,\beta,\gamma}^{(L)}
=\tr(\chi_{\alpha,\beta,\gamma}^L)$.
Each weighted copy of a target contributes its weight to the power $L$;
weights are summed when they multiply the same target. This is the
power-sum mechanism of
Theorem~\ref{thm:asym_power_sum_coefficients}.
The results below establish the tensor decompositions and periodic trace
identities for these examples, rather than a classification of
length-dependent RFPs.

\begin{definition}[The idempotents of an adapted basis]
	\label{def:asymex_basis_proj}
	\lean{MPSTensor.basisProj, P6Compression.pairReorder, P6Compression.pairU,
		P6Compression.pairUinv, P6Compression.pairProjInt,
		P6Compression.pairBlockInt}
	\leanok
	\uses{prop:asym_projector_weighted_sum}
	For mutually inverse matrices $U$ and $U^{-1}$, define
	$\Pi_j=U^{-1}E_{jj}U$. This is the product of the $j$-th column of
	$U^{-1}$ with the $j$-th row of $U$.
	In both examples below, a stacked bond coordinate is a quadruple of
	bits. Reorder it into the inner bond formed by the two middle bits
	and the outer bond formed by the two outer bits. Choose the adapted
	basis of the inner bond whose first two vectors are
	$(1,0,0,\pm1)$ and whose last two vectors are the remaining coordinate
	vectors. The two idempotents associated with the Bell vectors carry
	the weighted target tensors. Tensor products of these idempotents
	with the target matrices are then expressed in the original stacked
	bond coordinates.
\end{definition}

\begin{lemma}[The idempotents of an adapted basis are mutually orthogonal]
	\label{lem:asymex_basis_proj}
	\lean{MPSTensor.basisProj_mul_basisProj_self,
		MPSTensor.basisProj_mul_basisProj_of_ne, MPSTensor.mul_basisProj_mul}
	\leanok
	\uses{def:asymex_basis_proj}
	The idempotents satisfy $\Pi_j^2=\Pi_j$,
	$\Pi_j\Pi_k=0$ for $j\ne k$, and $U\Pi_jU^{-1}=E_{jj}$.
\end{lemma}

\begin{proof}
	\leanok
	The idempotent is $U^{-1}E_{jj}U$, so conjugation by $U$ returns
	$E_{jj}$, and the products reduce to $E_{jj}E_{kk}$, which is $E_{jj}$
	for $j=k$ and zero otherwise.
\end{proof}

\begin{theorem}[From an integer matrix identity to a weighted direct sum]
	\label{thm:asymex_stacked_pair}
	\lean{P6Compression.stackedPair_letter_identity}
	\leanok
	\uses{def:asymex_basis_proj, prop:asym_projector_weighted_sum}
	Suppose the stacked tensor is $c_B$ times an integer tensor
	$\widehat B$, and each target tensor $C_s$ is $c_C$ times an integer
	tensor $\widehat C_s$, with $s\in\{0,1\}$. Suppose there are a nonzero
	integer $n$ and integers $k_s$ such that, for every $s$ and every
	physical index $a$,
	\begin{align}
		n\,\mu_s\,\frac{c_C}{2} & =k_s\,c_B,
		\label{eq:asymex_stacked_pair_scalar}                            \\
		n\,\widehat B^a         & =\sum_sk_s\,\bigl(\widehat\Pi_s\otimes
		\widehat C_s^a\bigr),
		\label{eq:asymex_stacked_pair_letters}
	\end{align}
	where $\widehat\Pi_s=2\Pi_s$ and the right-hand side of
	(\ref{eq:asymex_stacked_pair_letters}) is expressed in the stacked
	bond coordinates. Then
	$B^a=\sum_s\mu_s(\Pi_s\otimes C_s^a)$ in those coordinates, which is
	the hypothesis of Proposition~\ref{prop:asym_projector_weighted_sum}.
\end{theorem}

\begin{proof}
	\leanok
	\uses{lem:asymex_basis_proj}
	Expand the right-hand side entrywise using the idempotents of
	Lemma~\ref{lem:asymex_basis_proj}. Since
	$\Pi_s=\widehat\Pi_s/2$ and $C_s^a=c_C\widehat C_s^a$,
	its $s$-th summand is
	$\mu_sc_C(\widehat\Pi_s\otimes\widehat C_s^a)/2$.
	Multiplying by $n$ and applying
	(\ref{eq:asymex_stacked_pair_scalar}) and
	(\ref{eq:asymex_stacked_pair_letters}) gives
	$c_Bn\widehat B^a=nB^a$. Cancel the nonzero scalar $n$.
\end{proof}

\begin{theorem}[Normality from scaled matrix units]
	\label{thm:asymex_normal_units}
	\lean{P6Compression.isNormal_of_single_eq_smul}
	\leanok
	\uses{def:normal}
	Let a tensor be one nonzero scalar times an integer tensor, and suppose
	that for every position $(x,y)$ of its bond algebra some matrix of the
	integer tensor is a nonzero integer multiple of $E_{xy}$. Then the
	one-site matrices of the original tensor span the full matrix algebra
	of the bond space, so the tensor is normal at blocking length one in
	the sense of Definition~\ref{def:normal}.
	This conclusion does not impose spectral-radius normalization;
	the term NT in the RFP canonical form additionally requires that
	normalization.
\end{theorem}

\begin{proof}
	\leanok
	Dividing the designated letter by its nonzero scalar factor exhibits
	each matrix unit in the span of the one-site matrices, and the matrix
	units span the whole algebra.
\end{proof}

\section{The one-label fixed point}%
\label{sec:asymex_one_label}

The first example has four one-site matrices:
\begin{align}
	A^0 & =\tfrac45E_{00},
	\label{eq:asymex_one_label_matrices} \\
	A^1 & =\tfrac45E_{11}, \\
	A^2 & =\tfrac35E_{01}, \\
	A^3 & =\tfrac35E_{10}.
\end{align}
The squared scales are $(1\pm\lambda)/2$ for $\lambda=7/25$.
Reading the tensor in the vertical direction gives
$M^{(i,j)}=A^i\otimes A^j$, with sixteen physical index pairs.
Its stacked product factorizes as $A^i\otimes T\otimes A^k$, where
$T=\sum_jA^j\otimes A^j$ acts on the inner bond and has spectrum
$\{1,7/25,0,0\}$. Here $T$ denotes this matrix, not a tpCPM in the
definition of an RFP. Its spectral decomposition gives two weighted
copies of the same target tensor, with weights $1$ and $7/25$.
Thus the associated periodic structure coefficient is
$c^{(L)}=\tr(\chi^L)=1+(7/25)^L$ with
$\chi=\diag(1,7/25)$, rather than a sum of contributions from
inequivalent target tensors.

\begin{definition}[The one-label tensors]
	\label{def:asymex_one_label}
	\lean{P6Compression.oneLabelM, P6Compression.oneLabelTarget,
		P6Compression.oneLabelStacked, P6Compression.oneLabelBlocks,
		P6Compression.oneLabelWeights}
	\leanok
	\uses{def:mpo_tensor, def:mpo_to_mps, def:mpdo_operator_product}
	The tensor read in the vertical direction is
	$M^{(i,j)}=A^i\otimes A^j$, with the matrices
	(\ref{eq:asymex_one_label_matrices}), of bond dimension four.
	Viewed as a matrix product tensor over the sixteen physical index
	pairs, it is the target tensor. Let $B$ be its stacked product with
	itself, of bond dimension sixteen. The two target tensors are copies
	of $M$ with weights $\mu_0=1$ and $\mu_1=7/25$.
\end{definition}

\begin{center}
	% Formula: the periodic contraction of the vertical tensor $M$ with
	% the inner-bond state $\sigma=\sum_b\lambda_b\Pi_b$ inserted on
	% every bond, whose weights are the spectrum $\{1,7/25,0,0\}$ of
	% $T=\sum_jA^j\otimes A^j$ (\S\ref{sec:asymex_rfp_basis}); this is
	% the spectral mechanism behind
	% Theorem~\ref{thm:asymex_one_label_compression}.
	% Source: \S\ref{sec:asymex_one_label}, introductory paragraph.
	% Ink-to-index map: the two explicit dot nodes are the vertical
	% tensor $M$, each carrying its physical leg of the pair alphabet;
	% each ring is the inner-bond state $\sigma$, with no physical leg
	% of its own; the trailing dots abbreviate the remaining alternating
	% $M$-$\sigma$ pairs of the periodic ring.
	% Contraction set: every virtual bond, alternating $M$ and $\sigma$,
	% closed around the ring; no leg is left open.
	% Boundary signature: (virtual-west=0 | virtual-east=0 |
	% physical-up=2 | physical-down=0) for the finite schematic shown.
	\tenkzkernel
	\begin{tenkz}[cols=6, boundary=periodic, physical=up]
		\tn[label pos=135, ports={90:physical:$a_1$}]{M} &
		\tn[skin=ring]{\sigma} &
		\tn[label pos=135, ports={90:physical:$a_2$}]{M} &
		\tn[skin=ring]{\sigma} &
		\tn[skin=dots]{} & \tn[skin=ring]{\sigma}
	\end{tenkz}
	\qquad $\sigma=\sum_b\lambda_b\,\Pi_b$
\end{center}

\begin{theorem}[Normality of the one-label target]
	\label{thm:asymex_one_label_normal}
	\lean{P6Compression.oneLabelTarget_isNormal}
	\leanok
	\uses{def:asymex_one_label, def:normal}
	The vertical tensor is normal at blocking length one in the sense of
	Definition~\ref{def:normal}: its sixteen matrices are scaled matrix
	units covering every position of the four-by-four matrix algebra.
\end{theorem}

\begin{proof}
	\leanok
	\uses{thm:asymex_normal_units}
	Each of the sixteen positions is realized by one of the sixteen letters
	with a nonzero integer factor, a finite verification over integer
	matrices; Theorem~\ref{thm:asymex_normal_units} then gives normality.
\end{proof}

\begin{theorem}[Split compression of the one-label fixed point]
	\label{thm:asymex_one_label_compression}
	% Principal declaration: P6Compression.oneLabel_isReduction.
	\lean{P6Compression.oneLabelCompression, P6Compression.oneLabel_isReduction,
		P6Compression.oneLabel_left_mul_right_of_ne, P6Compression.oneLabel_z_eq,
		P6Compression.oneLabel_dim_eq, P6Compression.oneLabel_remainder,
		P6Compression.oneLabel_mul_right_eq_right_mul,
		P6Compression.oneLabel_left_mul_eq_mul_left}
	\leanok
	\uses{def:asymex_one_label, thm:asym_multiblock}
	The stacked tensor $B$ of Definition~\ref{def:asymex_one_label} has the
	block triangular form of Theorem~\ref{thm:asym_multiblock} for the two
	weighted copies $C_s=\mu_sM$, with an eight-dimensional zero
	complement, $z=8$, and $16=4+4+8$. The compression pairs are mutually
	biorthogonal and satisfy $W_sB^wV_s=C_s^w$ for every word.
	The remainder vanishes identically, and the compression maps
	intertwine the one-site matrices in both directions:
	$B^aV_s=V_s(\mu_sM^a)$ and $W_sB^a=(\mu_sM^a)W_s$.
\end{theorem}

\begin{proof}
	\leanok
	\uses{thm:asymex_stacked_pair, prop:asym_projector_weighted_sum}
	Multiplying the stacked tensor by two and expanding it in the adapted
	basis of Section~\ref{sec:asymex_rfp_basis} gives the integer matrix
	identity of Theorem~\ref{thm:asymex_stacked_pair} with the integer
	coefficients $25$ and $7$, a finite verification over the sixteen
	physical indices. The matching scalar identity holds because the two
	weights are $1$ and $7/25$.
	Proposition~\ref{prop:asym_projector_weighted_sum} then supplies the
	compression maps, the vanishing remainder and the one-site
	intertwining identities, with a zero complement of dimension
	$8=2\cdot4$.
\end{proof}

\begin{center}
	% Formula: B^aV_s=V_s(\mu_sM^a) and W_sB^a=(\mu_sM^a)W_s of
	% Theorem~\ref{thm:asymex_one_label_compression}, for the weights
	% $\mu_s\in\{1,7/25\}$.
	% Source: Theorem~\ref{thm:asymex_one_label_compression}.
	% Ink-to-index map: the ring is the compression map $V_s$ or $W_s$;
	% the square node is the stacked source $B$ or the weighted target
	% $\mu_sM$.
	% Contraction set: the single virtual bond joining the ring to the
	% square node in each pair; the physical pair leg $a$ remains open.
	% Boundary signature: (virtual-west=1 | virtual-east=1 |
	% physical-up=1 | physical-down=0) for each pictured equality.
	\tenkzkernel
	\begin{tenkz}[rows={wire}, cols=2, west=open, east=open]
		\tn[label pos=s, ports={90:physical:$a$}]{B} & \tn[skin=ring]{V_s}
	\end{tenkz}
	\quad $=$ \quad
	\begin{tenkz}[rows={wire}, cols=2, west=open, east=open]
		\tn[skin=ring]{V_s} & \tn[label pos=s, ports={90:physical:$a$}]{\mu_sM}
	\end{tenkz}
	\\[2ex]
	\begin{tenkz}[rows={wire}, cols=2, west=open, east=open]
		\tn[skin=ring]{W_s} & \tn[label pos=s, ports={90:physical:$a$}]{B}
	\end{tenkz}
	\quad $=$ \quad
	\begin{tenkz}[rows={wire}, cols=2, west=open, east=open]
		\tn[label pos=s, ports={90:physical:$a$}]{\mu_sM} & \tn[skin=ring]{W_s}
	\end{tenkz}
\end{center}

\begin{theorem}[The one-label structure coefficient]
	\label{thm:asymex_one_label_trace}
	\lean{P6Compression.oneLabel_trace_evalWord}
	\leanok
	\uses{def:asymex_one_label}
	For every nonempty word $w$ over the sixteen-letter alphabet,
	\begin{align}
		\tr(B^w) & =\left(1+(\tfrac{7}{25})^{|w|}\right)\tr(M^w).
		\label{eq:asymex_one_label_trace}
	\end{align}
\end{theorem}

\begin{proof}
	\leanok
	\uses{thm:asymex_one_label_compression, prop:asym_compression_trace}
	Proposition~\ref{prop:asym_compression_trace} applied to the compression
	in Theorem~\ref{thm:asymex_one_label_compression} gives the sum of the
	two weighted target traces. Pulling each weight out of its word
	evaluation contributes $1$ and $(7/25)^{|w|}$.
\end{proof}

\section{The graded quantum-dimer twist}%
\label{sec:asymex_dimer}

In the two-sector graded quantum-dimer twist, a horizontal bond index is a
triple $(p,p',k)$ of bits: the ket and bra indices of the left half-bond
and a sector label. The vertical site space is $\C^2\otimes\C^2$.
Each nonzero matrix of the sector-$f$ tensor is a scaled matrix unit
with scale $\tfrac12C_k[p][p']\tau_k[f]$, where
$C_0=xP_++yP_-$, $C_1=xP_+-yP_-$,
$\tau_0=(1,1)$ and $\tau_1=(1,-1)$.
Matrices whose two horizontal bond indices have different sector labels
vanish.

The target tensors for the stacked product of sectors $f$ and $f'$ are
the sectors $f+f'$ and $f+f'+1$, with addition modulo two, and their
weights are $x/2$ and $y/2$. At $x=7/8$, $y=1/8$, these weights are
$7/16$ and $1/16$. Unlike the one-label example, the two contributions
multiply different target tensors. For the product of sector zero with
itself, the periodic structure coefficients are
$c_{0,0,0}^{(L)}=(7/16)^L$ and
$c_{0,0,1}^{(L)}=(1/16)^L$, the power sums of the one-dimensional
positive diagonal matrices $\chi_{0,0,0}=(7/16)$ and
$\chi_{0,0,1}=(1/16)$, respectively.
The compression and trace theorems below concern this pair of sectors.

\begin{definition}[The sector tensors and their stacked product]
	\label{def:asymex_dimer}
	\lean{P6Compression.dimerM, P6Compression.dimerTarget,
		P6Compression.dimerStacked, P6Compression.dimerBlocks,
		P6Compression.dimerWeights}
	\leanok
	\uses{def:mpo_tensor, def:mpo_to_mps, def:mpdo_operator_product}
	At $x=7/8$ and $y=1/8$, the sector-$f$ tensor has, at the index pair
	$(a,b)$ of horizontal bonds, the scaled matrix unit whose row is the
	pair of ket bits of $a$ and $b$, whose column is the pair of bra bits,
	and whose scale is $\tfrac12C_k[p][p']\tau_k[f]$ for their common
	sector label $k$. It vanishes when the two sector labels differ.
	Read in the vertical direction, over the sixty-four index pairs,
	this gives the target tensor $M_f$.
	The stacked product of the sector-$f$ and sector-$f'$ tensors has
	bond dimension sixteen. The two target tensors are $M_{f+f'}$ and
	$M_{f+f'+1}$ with weights $7/16$ and $1/16$, respectively, where
	sector addition is modulo two.
\end{definition}

\begin{center}
	% Formula: the periodic contraction of the sector-$f$ tensor $M_f$
	% with the inner-bond state $\sigma=\sum_b\lambda_b\Pi_b$ inserted
	% on every bond, the same adapted-basis mechanism of
	% \S\ref{sec:asymex_rfp_basis} that gives the weights $7/16$ and
	% $1/16$.
	% Source: \S\ref{sec:asymex_dimer}, introductory paragraph.
	% Ink-to-index map: the two explicit dot nodes are the sector tensor
	% $M_f$, each carrying its physical leg of the pair alphabet; each
	% ring is the inner-bond state $\sigma$, with no physical leg of its
	% own; the trailing dots abbreviate the remaining alternating
	% $M_f$-$\sigma$ pairs of the periodic ring.
	% Contraction set: every virtual bond, alternating $M_f$ and
	% $\sigma$, closed around the ring; no leg is left open.
	% Boundary signature: (virtual-west=0 | virtual-east=0 |
	% physical-up=2 | physical-down=0) for the finite schematic shown.
	\tenkzkernel
	\begin{tenkz}[cols=6, boundary=periodic, physical=up]
		\tn[label pos=135, ports={90:physical:$a_1$}]{M_f} &
		\tn[skin=ring]{\sigma} &
		\tn[label pos=135, ports={90:physical:$a_2$}]{M_f} &
		\tn[skin=ring]{\sigma} &
		\tn[skin=dots]{} & \tn[skin=ring]{\sigma}
	\end{tenkz}
	\qquad $\sigma=\sum_b\lambda_b\,\Pi_b$
\end{center}

\begin{theorem}[Normality of the sectors]
	\label{thm:asymex_dimer_normal}
	\lean{P6Compression.dimerTarget_isNormal}
	\leanok
	\uses{def:asymex_dimer, def:normal}
	Each sector tensor is normal at blocking length one in the sense of
	Definition~\ref{def:normal}: its thirty-two nonzero matrices are
	scaled matrix units covering every position of the four-by-four
	matrix algebra.
\end{theorem}

\begin{proof}
	\leanok
	\uses{thm:asymex_normal_units}
	Sixteen letters lying in the sector label zero, where the scale does
	not depend on $f$, realize the sixteen positions with nonzero integer
	factors, a finite verification over integer matrices;
	Theorem~\ref{thm:asymex_normal_units} then gives normality.
\end{proof}

\begin{theorem}[Split compression of the dimer fusion]
	\label{thm:asymex_dimer_compression}
	% Principal declaration: P6Compression.dimer_isReduction.
	\lean{P6Compression.dimerCompression, P6Compression.dimer_isReduction,
		P6Compression.dimer_left_mul_right_of_ne, P6Compression.dimer_z_eq,
		P6Compression.dimer_dim_eq, P6Compression.dimer_remainder,
		P6Compression.dimer_mul_right_eq_right_mul,
		P6Compression.dimer_left_mul_eq_mul_left}
	\leanok
	\uses{def:asymex_dimer, thm:asym_multiblock}
	For the product of sector zero with itself, let $B$ be the stacked
	tensor of Definition~\ref{def:asymex_dimer}.
	It has the block triangular form of
	Theorem~\ref{thm:asym_multiblock} for the two weighted target tensors
	$C_s=\mu_sM_s$, where $\mu_0=7/16$ and $\mu_1=1/16$, with an
	eight-dimensional zero complement, $z=8$, and $16=4+4+8$.
	The compression pairs are mutually biorthogonal and satisfy
	$W_sB^wV_s=C_s^w$ for every word. The remainder vanishes identically,
	and the compression maps intertwine the one-site matrices in both
	directions.
\end{theorem}

\begin{proof}
	\leanok
	\uses{thm:asymex_stacked_pair, prop:asym_projector_weighted_sum}
	A matrix of the stacked tensor vanishes unless its two outer
	horizontal bond indices carry the same sector label. For a matching
	pair, the sum over the intermediate index runs over the four bond
	indices of that label. Multiplying the resulting matrix by two and
	expanding it in the adapted basis of
	Section~\ref{sec:asymex_rfp_basis} gives the integer matrix identity
	of Theorem~\ref{thm:asymex_stacked_pair} with coefficients $7$ and
	$1$, a finite verification over the sixty-four physical indices in
	each of the two sector-label cases. The matching scalar identity
	holds because the weights are $7/16$ and $1/16$.
	Proposition~\ref{prop:asym_projector_weighted_sum} then supplies the
	compression maps, the vanishing remainder and the one-site
	intertwining identities.
\end{proof}

\begin{center}
	% Formula: the sitewise intertwiners B^aV_s=V_s(\mu_sM_s^a) and
	% W_sB^a=(\mu_sM_s^a)W_s of Theorem~\ref{thm:asymex_dimer_compression},
	% for the two weighted target sectors $s\in\{f+f',f+f'+1\}$ with
	% weights $\mu_s\in\{7/16,1/16\}$.
	% Source: Theorem~\ref{thm:asymex_dimer_compression}.
	% Ink-to-index map: the ring is the compression map $V_s$ or $W_s$;
	% the square node is the stacked source $B$ or the weighted sector
	% target $\mu_sM_s$.
	% Contraction set: the single virtual bond joining the ring to the
	% square node in each pair; the physical pair leg $a$ remains open.
	% Boundary signature: (virtual-west=1 | virtual-east=1 |
	% physical-up=1 | physical-down=0) for each pictured equality.
	\tenkzkernel
	\begin{tenkz}[rows={wire}, cols=2, west=open, east=open]
		\tn[label pos=s, ports={90:physical:$a$}]{B} & \tn[skin=ring]{V_s}
	\end{tenkz}
	\quad $=$ \quad
	\begin{tenkz}[rows={wire}, cols=2, west=open, east=open]
		\tn[skin=ring]{V_s} &
		\tn[label pos=s, ports={90:physical:$a$}]{\mu_sM_s}
	\end{tenkz}
	\\[2ex]
	\begin{tenkz}[rows={wire}, cols=2, west=open, east=open]
		\tn[skin=ring]{W_s} & \tn[label pos=s, ports={90:physical:$a$}]{B}
	\end{tenkz}
	\quad $=$ \quad
	\begin{tenkz}[rows={wire}, cols=2, west=open, east=open]
		\tn[label pos=s, ports={90:physical:$a$}]{\mu_sM_s} & \tn[skin=ring]{W_s}
	\end{tenkz}
\end{center}

\begin{theorem}[The dimer structure coefficients]
	\label{thm:asymex_dimer_trace}
	\lean{P6Compression.dimer_trace_evalWord}
	\leanok
	\uses{def:asymex_dimer}
	For the product of sector zero with itself and every nonempty word
	$w$ over the sixty-four-letter alphabet,
	\begin{align}
		\tr(B^w)
		 & =(\tfrac{7}{16})^{|w|}\tr(M_0^w)
		+(\tfrac{1}{16})^{|w|}\tr(M_1^w),
		\label{eq:asymex_dimer_trace}
	\end{align}
	where $M_0$ and $M_1$ are the two sector tensors.
\end{theorem}

\begin{proof}
	\leanok
	\uses{thm:asymex_dimer_compression, prop:asym_compression_trace}
	Proposition~\ref{prop:asym_compression_trace} applied to the compression
	in Theorem~\ref{thm:asymex_dimer_compression} gives the sum of the two
	weighted target traces. Pulling each weight out of its word
	evaluation gives the factors $(7/16)^{|w|}$ and $(1/16)^{|w|}$ in
	(\ref{eq:asymex_dimer_trace}).
\end{proof}

\begin{remark}[Scope of the dimer fusion identities]
	\label{rem:asymex_dimer_scope}
	Theorems~\ref{thm:asymex_dimer_compression}
	and~\ref{thm:asymex_dimer_trace} concern the product of sector zero
	with itself. The remaining three pairs of sectors satisfy the same
	one-site matrix identity, with the same adapted basis and weights,
	and with targets $M_{f+f'}$ and $M_{f+f'+1}$. For each pair this
	requires a separate verification over the sixty-four physical
	indices, which is not carried out here.
\end{remark}
%

\section{Structure coefficients and boundary matrices}
\label{sec:asymex_coefficients_boundaries}

\subsection{Representative rescaling}

We use the structure coefficients
$c_{\alpha,\beta,\gamma}^{(L)}=\tr(\chi_{\alpha,\beta,\gamma}^{L})$
and trace scalars $m_\alpha$ of
\cite[Theorem~4.14(ii)]{Cirac2016MPDO_arXiv}.
The following identities concern these coefficients, not the construction
of a renormalization fixed point.

\begin{definition}[Rescaling the structure coefficients]
    \label{def:asymex_rescaling}
    \lean{MPOTensor.DiagonalChiFamily.rescale,
        MPOTensor.DiagonalChiFamily.rescale_dim,
        MPOTensor.DiagonalChiFamily.rescale_entry,
        MPOTensor.BNTLabelCoefficientFamily.rescale,
        MPOTensor.BNTLabelCoefficientFamily.rescale_coeff,
        MPOTensor.BNTLabelTraceScalarFamily.rescale,
        MPOTensor.BNTLabelTraceScalarFamily.rescale_traceScalar}
    \leanok
    \uses{def:mpdo_diagonal_chi_family,
        def:mpdo_bnt_label_coefficient_family,
        def:mpdo_bnt_label_trace_scalar_family}
    Let $\Lambda$ be a label set and let $s_\alpha\in\C$
    for every $\alpha\in\Lambda$. In the algebraic formulas below,
    division is extended by $z/0=0$. Their interpretation as changes
    of tensor representatives requires all scales to be nonzero.
    Rescale the diagonal entries, coefficients, and trace scalars by
    \begin{align}
        \chi'_{\alpha,\beta,\gamma,k}
                  & =\frac{s_\alpha s_\beta}{s_\gamma}
        \chi_{\alpha,\beta,\gamma,k},
        \label{eq:asymex_chi_rescale}                                 \\
        c'{}_{\alpha,\beta,\gamma}^{(L)}
                  & =\left(\frac{s_\alpha s_\beta}{s_\gamma}\right)^L
        c_{\alpha,\beta,\gamma}^{(L)},
        \label{eq:asymex_c_rescale}                                   \\
        m'_\alpha & =\frac{m_\alpha}{s_\alpha}.
        \label{eq:asymex_m_rescale}
    \end{align}
    The multiplicity $r_{\alpha,\beta,\gamma}$ is unchanged, and
    $0\leq k<r_{\alpha,\beta,\gamma}$.
\end{definition}

\begin{theorem}[Trace powers and positivity under rescaling]
    \label{thm:asymex_rescale_trace}
    \lean{MPOTensor.DiagonalChiFamily.tracePowerCoeff_rescale,
        MPOTensor.BNTLabelCoefficientFamily.ofChi_rescale,
        MPOTensor.DiagonalChiFamily.PosEntries.rescale}
    \leanok
    \uses{def:asymex_rescaling}
    Taking trace-power coefficients commutes with rescaling: for every
    $L\in\N$, including $L=0$,
    \begin{align}
        \tr((\chi'_{\alpha,\beta,\gamma})^L)
         & =
        \left(\frac{s_\alpha s_\beta}{s_\gamma}\right)^L
        \tr(\chi_{\alpha,\beta,\gamma}^L).
    \end{align}
    If the original diagonal entries are positive real numbers and
    every $s_\alpha$ is positive real, the rescaled entries are positive.
\end{theorem}
\begin{proof}
    \leanok
    \uses{def:asymex_rescaling}
    Expand the trace as the sum of the $L$th powers of the diagonal
    entries. Each summand acquires the factor
    $(s_\alpha s_\beta/s_\gamma)^L$, which factors out of the finite
    sum. For positive real scales, the ratio
    $s_\alpha s_\beta/s_\gamma$ is positive, so multiplication by it
    preserves positivity.
\end{proof}

\begin{theorem}[Preservation of the idempotent coefficient equation]
    \label{thm:asymex_rescale_idempotent}
    \lean{MPOTensor.BNTLabelTraceScalarFamily.hasIdempotentCoefficientForm_rescale_iff,
        MPOTensor.BNTLabelTraceScalarFamily.rescale_coefficient_sum}
    \leanok
    \uses{def:asymex_rescaling,
        def:mpdo_bnt_label_idempotent_coefficient_form}
    Suppose $\Lambda$ is finite and every $s_\alpha$ is nonzero.
    The vector $m'$ is idempotent for $c'{}^{(1)}$ if and only if
    $m$ is idempotent for $c^{(1)}$. Explicitly, the equations
    \begin{align}
        m'_\gamma
         & =\sum_{\alpha,\beta}
        c'{}_{\alpha,\beta,\gamma}^{(1)}m'_\alpha m'_\beta,
        \\
        m_\gamma
         & =\sum_{\alpha,\beta}
        c_{\alpha,\beta,\gamma}^{(1)}m_\alpha m_\beta
    \end{align}
    are equivalent when imposed for every $\gamma\in\Lambda$.
\end{theorem}
\begin{proof}
    \leanok
    \uses{def:asymex_rescaling,
        def:mpdo_bnt_label_idempotent_coefficient_form}
    Substituting~(\ref{eq:asymex_c_rescale})
    and~(\ref{eq:asymex_m_rescale}) and cancelling the nonzero
    scales in each summand gives
    \begin{align}
        \sum_{\alpha,\beta}
        c'{}_{\alpha,\beta,\gamma}^{(1)}m'_\alpha m'_\beta
         & =
        \frac{1}{s_\gamma}
        \sum_{\alpha,\beta}
        c_{\alpha,\beta,\gamma}^{(1)}m_\alpha m_\beta.
    \end{align}
    The left side of the idempotent equation is
    $m'_\gamma=m_\gamma/s_\gamma$. Cancellation of $s_\gamma$
    proves both implications.
\end{proof}

Multiplying a normalized normal-tensor representative by a positive
scalar different from one leaves the spectral-radius-one convention.
These coefficient identities neither construct normalized
representatives nor establish a new renormalization fixed point.

\subsection{Coefficients from the Ising fusion rules}

\begin{definition}[Structure coefficients from the Ising fusion rules]
    \label{def:asymex_ising_coefficients}
    \lean{TNLean.Algebra.IsingCrossing.Label,
        TNLean.Algebra.IsingCrossing.fusionMultiplicity,
        TNLean.Algebra.IsingCrossing.coefficient}
    \leanok
    Let the labels be $1,\psi,\sigma$, and let $N_{ab}^{c}\in\N$
    be the multiplicity of $c$ in the Ising fusion rules. For each
    label $a$, these rules are
    \begin{align}
        1\times a          & =a\times 1=a,             \\
        \psi\times\psi     & =1,                       \\
        \psi\times\sigma   & =\sigma\times\psi=\sigma, \\
        \sigma\times\sigma & =1+\psi.
    \end{align}
    For arbitrary weights $w_a\in\C$ and $L\in\N$, define
    \begin{align}
        c_{a,b,c}^{(L)}
         & =\sum_r\sum_t w_r^L N_{ar}^{t}N_{tb}^{c}.
        \label{eq:asymex_ising_coeff}
    \end{align}
    Both sums run over $\{1,\psi,\sigma\}$.
    The definition includes $L=0$ and imposes no positivity or
    nonvanishing condition on the weights.
\end{definition}

\begin{theorem}[The three outputs for two sigma labels]
    \label{thm:asymex_ising_sigma}
    \lean{TNLean.Algebra.IsingCrossing.coefficient_sigma_sigma_one,
        TNLean.Algebra.IsingCrossing.coefficient_sigma_sigma_psi,
        TNLean.Algebra.IsingCrossing.coefficient_sigma_sigma_sigma,
        TNLean.Algebra.IsingCrossing.fusionMultiplicity_sigma_sigma_sigma}
    \leanok
    \uses{def:asymex_ising_coefficients}
    For every choice of complex weights and every $L\in\N$,
    \begin{align}
        c_{\sigma,\sigma,1}^{(L)}      & =w_1^L+w_\psi^L, \\
        c_{\sigma,\sigma,\psi}^{(L)}   & =w_1^L+w_\psi^L, \\
        c_{\sigma,\sigma,\sigma}^{(L)} & =2w_\sigma^L.
    \end{align}
    In contrast, the base fusion multiplicity is
    $N_{\sigma\sigma}^{\sigma}=0$.
\end{theorem}
\begin{proof}
    \leanok
    \uses{def:asymex_ising_coefficients}
    Expand both three-element sums in~(\ref{eq:asymex_ising_coeff})
    and substitute the fusion table. For output $1$ or $\psi$, the
    surviving inserted labels are $1$ and $\psi$. For output $\sigma$,
    the inserted label $\sigma$ contributes through both intermediate
    labels $1$ and $\psi$.
\end{proof}

\begin{theorem}[Associativity of the structure coefficients]
    \label{thm:asymex_ising_assoc}
    \lean{TNLean.Algebra.IsingCrossing.coefficient_assoc}
    \leanok
    \uses{def:asymex_ising_coefficients}
    For arbitrary complex weights, every $L\in\N$, and all labels
    $a,b,c,d$,
    \begin{align}
        \sum_e c_{a,b,e}^{(L)}c_{e,c,d}^{(L)}
         & =\sum_f c_{b,c,f}^{(L)}c_{a,f,d}^{(L)}.
    \end{align}
\end{theorem}
\begin{proof}
    \leanok
    \uses{def:asymex_ising_coefficients}
    Consider the $3^4$ choices of $a,b,c,d$. In each case, expand
    the finite sums and substitute the fusion multiplicities.
    The resulting expressions agree as polynomials in
    $w_1^L,w_\psi^L,w_\sigma^L$.
\end{proof}

These examples use the structure-coefficient notation of
\cite{Cirac2016MPDO_arXiv}. They do not give a tensor realization,
comparison matrices, or a classification for fixed $F$-symbols.

\subsection{Explicit scalar three-cocycles on the group of order two}

We use the three-cocycle and fusion-gauge conventions of
\cite{GarreRubioSchuch2024DomainWall}. These concern scalar associators,
rather than the structure coefficients of the preceding subsections.

\begin{definition}[Two scalar three-cochains]
    \label{def:asymex_z2_cocycles}
    \lean{TNLean.Algebra.ScalarThreeCochain.cyclicTwoCocycle,
        TNLean.Algebra.ScalarThreeCochain.cyclicTwoCocycle_generator,
        TNLean.Algebra.ScalarThreeCochain.cyclicTwoCocycle_zero}
    \leanok
    \uses{def:mpug_three_cocycle}
    Write $\Z_2=\{e,s\}$ multiplicatively, with $s^2=e$.
    Identify $e,s$ with the bits $0,1$ only in exponents.
    For $p\in\{0,1\}$, define
    \begin{align}
        \omega_p(a,b,c) & =(-1)^{pabc}.
    \end{align}
    Thus $\omega_p(s,s,s)=(-1)^p$, all other values are one,
    and $\omega_0$ is the constant cochain one.
\end{definition}

\begin{theorem}[Normalization, cocycle equation, and unit modulus]
    \label{thm:asymex_z2_properties}
    \lean{TNLean.Algebra.ScalarThreeCochain.cyclicTwoCocycle_isNormalized,
        TNLean.Algebra.ScalarThreeCochain.cyclicTwoCocycle_isCocycle,
        TNLean.Algebra.ScalarThreeCochain.cyclicTwoCocycle_norm}
    \leanok
    \uses{def:asymex_z2_cocycles}
    Each $\omega_p$ is normalized, has unit modulus, and satisfies,
    for all $a,b,c,d\in\Z_2$,
    \begin{align}
        \omega_p(a,b,c)\omega_p(a,bc,d)\omega_p(b,c,d)
         & =\omega_p(ab,c,d)\omega_p(a,b,cd).
    \end{align}
    Normalization means that its value is one whenever any argument is $e$.
\end{theorem}
\begin{proof}
    \leanok
    \uses{def:asymex_z2_cocycles}
    An identity argument excludes the triple $(s,s,s)$, proving
    normalization. For the cocycle equation, check the two choices
    for each of $a,b,c,d$ and for $p$, using $s^2=e$.
    Every value is either one or a power of $-1$, so its modulus is one.
\end{proof}

\begin{theorem}[Distinct scalar gauge classes]
    \label{thm:asymex_z2_classes}
    \lean{TNLean.Algebra.ScalarThreeCochain.cyclicTwoCocycle_cohomologousTo_iff,
        TNLean.Algebra.ScalarThreeCochain.cyclicTwoCocycle_one_not_isTrivialGaugeClass,
        TNLean.Algebra.ScalarThreeCochain.cyclicTwoCocycle_sigma_generator,
        TNLean.Algebra.ScalarThreeCochain.cyclicTwoCocycle_fusionGauge_generator}
    \leanok
    \uses{def:asymex_z2_cocycles,def:mpug_cohomologous_to,def:mpug_inversion_cochains}
    The cochains $\omega_p$ and $\omega_q$ are cohomologous if and
    only if $p=q$, even when the scalar two-cochain implementing the
    gauge change is not required to be normalized.
    The inversion scalar $\sigma_\omega(g)=\omega(g,g^{-1},g)$
    therefore satisfies $\sigma_{\omega_p}(s)=(-1)^p$.
    For a normalized scalar two-cochain $\beta$, the fusion-gauge
    transform satisfies
    \begin{align}
        (\beta\cdot\omega_p)(s,s,s) & =(-1)^p,
    \end{align}
    where the convention is
    \begin{align}
        (\beta\cdot\omega)(a,b,c)
         & =\frac{\beta(b,c)\beta(a,bc)}
        {\beta(a,b)\beta(ab,c)}\omega(a,b,c).
    \end{align}
\end{theorem}
\begin{proof}
    \leanok
    \uses{def:asymex_z2_cocycles,thm:asymex_z2_properties,
        def:mpug_cohomologous_to,thm:mpug_z2_trivial_scalar_class}
    Since $s^{-1}=s$, the inversion scalar equals
    $\omega_p(s,s,s)$. The criterion in Theorem~\ref{thm:mpug_z2_trivial_scalar_class}
    identifies triviality of the scalar gauge
    class with $\sigma_\omega(s)=1$. Hence $\omega_1$ is not
    gauge trivial. If $p\neq q$, one representative is $\omega_0=1$
    and the other is $\omega_1$; cohomology, using symmetry when
    needed, would contradict this nontriviality.
    If $p=q$, use reflexivity.
    Finally, substitute $s^2=e$ in the fusion-gauge coboundary
    factor and use $\beta(e,s)=\beta(s,e)=1$.
\end{proof}

No physical MPU or realization of an anomaly is constructed here.

\subsection{A Jordan extension detected by a boundary}

\begin{definition}[Jordan extension and split tensor]
    \label{def:asymex_jordan}
    \lean{MPSTensor.JordanBoundary.blocks,
        MPSTensor.JordanBoundary.tensor,
        MPSTensor.JordanBoundary.splitTensor,
        MPSTensor.JordanBoundary.boundary}
    \leanok
    Let $A=\{A^i\}_{i=0}^{d-1}$ be any tensor with $A^i\in\MN{D}$.
    In the standard two-block coordinates on $\C^{D+D}$, set
    \begin{align}
        B^i & =\begin{pmatrix}A^i&A^i\\0&A^i\end{pmatrix}, \\
        C^i & =\begin{pmatrix}A^i&0\\0&A^i\end{pmatrix},   \\
        Q   & =\begin{pmatrix}0&0\\\Id&0\end{pmatrix}.
    \end{align}
    Thus $C=A\oplus A$. For any word $w$, write $A^w$ for its
    ordered matrix product, with $A^\varnothing=\Id$.
\end{definition}

\begin{theorem}[Word products and boundary traces]
    \label{thm:asymex_jordan_traces}
    \lean{MPSTensor.JordanBoundary.evalWord_tensor,
        MPSTensor.JordanBoundary.evalWord_splitTensor,
        MPSTensor.JordanBoundary.trace_evalWord_tensor,
        MPSTensor.JordanBoundary.periodic_trace_eq,
        MPSTensor.JordanBoundary.boundary_trace_tensor,
        MPSTensor.JordanBoundary.boundary_trace_splitTensor}
    \leanok
    \uses{def:asymex_jordan}
    Every word $w$, including the empty word, satisfies
    \begin{align}
        B^w       & =\begin{pmatrix}A^w&|w|A^w\\0&A^w\end{pmatrix}, \\
        C^w       & =\begin{pmatrix}A^w&0\\0&A^w\end{pmatrix},      \\
        \tr(B^w)  & =\tr(C^w)=2\tr(A^w),                            \\
        \tr(QB^w) & =|w|\tr(A^w),                                   \\
        \tr(QC^w) & =0.
    \end{align}
\end{theorem}
\begin{proof}
    \leanok
    \uses{def:asymex_jordan}
    Induct on $w$. The empty word gives the identity matrix.
    Prepending a letter $i$ changes the upper-right block of $B^w$
    to $A^i(|w|A^w)+A^iA^w=(|w|+1)A^iA^w$.
    The split product retains zero off-diagonal blocks.
    Taking the sum of the diagonal-block traces proves the periodic
    formulas. Multiplication by $Q$ puts the upper-right word block
    in the lower-right block, giving the two boundary formulas.
\end{proof}

\begin{theorem}[Equal periodic traces and different boundary families]
    \label{thm:asymex_jordan_detection}
    \lean{MPSTensor.JordanBoundary.periodic_eq_boundary_ne,
        MPSTensor.JordanBoundary.boundary_trace_ne}
    \leanok
    \uses{def:asymex_jordan}
    Suppose there is a nonempty word $w$ with $\tr(A^w)\neq0$.
    Then $B$ and $C$ have equal periodic traces for every word,
    but the functions $v\mapsto\tr(QB^v)$ and
    $v\mapsto\tr(QC^v)$ differ.
    No normality assumption on $A$ is required.
\end{theorem}
\begin{proof}
    \leanok
    \uses{thm:asymex_jordan_traces}
    The periodic traces agree by
    Theorem~\ref{thm:asymex_jordan_traces}.
    At the specified word $w$, the boundary traces are
    $|w|\tr(A^w)\neq0$ and zero, respectively.
    Equality of the boundary functions would imply equality at $w$,
    a contradiction.
\end{proof}
%
